\documentclass[11pt,a4paper]{article}
\usepackage[T1]{fontenc}
\usepackage[utf8]{inputenc}
\usepackage{lmodern,microtype}
\usepackage[a4paper,margin=27mm,headheight=15pt]{geometry}
\usepackage{amsmath,amssymb,amsthm,mathtools,bm,mathrsfs}
\usepackage{booktabs,longtable,array,enumitem}
\usepackage[dvipsnames]{xcolor}
\usepackage{hyperref,bookmark,fancyhdr}
\hypersetup{colorlinks=true,linkcolor=MidnightBlue,citecolor=MidnightBlue,urlcolor=MidnightBlue,
 pdftitle={Surcomplex Analysis: Local Calculus, Coherent Hahn Functions, and Residue Theory},
 pdfauthor={OpenAI},pdfsubject={Analysis over the complexified surreal numbers}}
\definecolor{heading}{HTML}{163A50}
\usepackage{titlesec}
\titleformat{\section}{\Large\bfseries\color{heading}}{\thesection}{0.7em}{}
\titleformat{\subsection}{\large\bfseries\color{heading}}{\thesubsection}{0.7em}{}
\pagestyle{fancy}\fancyhf{}
\fancyhead[L]{\small Surcomplex analysis}\fancyhead[R]{\small Local and coherent theories}
\fancyfoot[C]{\thepage}
\setlength{\parskip}{4pt}\setlength{\parindent}{1.2em}
\setlist{itemsep=2pt,topsep=4pt}
\allowdisplaybreaks[2]
\newtheorem{theorem}{Theorem}[section]
\newtheorem{proposition}[theorem]{Proposition}
\newtheorem{lemma}[theorem]{Lemma}
\newtheorem{corollary}[theorem]{Corollary}
\theoremstyle{definition}\newtheorem{definition}[theorem]{Definition}
\newtheorem{example}[theorem]{Example}
\theoremstyle{remark}\newtheorem{remark}[theorem]{Remark}
\newcommand{\No}{\mathbf{No}}
\newcommand{\K}{\mathbb K}
\newcommand{\C}{\mathbb C}\newcommand{\R}{\mathbb R}\newcommand{\N}{\mathbb N}\newcommand{\Z}{\mathbb Z}
\newcommand{\A}{\mathscr H}\newcommand{\hol}{\mathcal O}
\newcommand{\fin}{\mathcal O_{\K}}\newcommand{\infinit}{\mathfrak m_{\K}}
\newcommand{\st}{\operatorname{st}}\newcommand{\supp}{\operatorname{supp}}
\newcommand{\Res}{\operatorname{Res}}\newcommand{\ord}{\operatorname{ord}}
\newcommand{\lc}{\operatorname{lc}}\newcommand{\id}{\operatorname{id}}
\newcommand{\ev}[1]{\widehat{#1}}
\newcommand{\hoint}{\mathop{\oint}\nolimits^{\!H}}
\newcommand{\hint}{\mathop{\int}\nolimits^{\!H}}
\newcommand{\abs}[1]{\left|#1\right|}
\newcommand{\set}[1]{\left\{#1\right\}}
\newcommand{\coeff}[2]{\left[#1\right]#2}
\newcommand{\vH}{v_H}
\newcommand{\halo}[1]{#1^{\mathrm h}}
\newcommand{\smallfield}{k_\Gamma}
\newcommand{\dd}{\,\mathrm d}
\numberwithin{equation}{section}

\begin{document}
\begin{titlepage}
\vspace*{17mm}
{\color{heading}\Huge\bfseries Surcomplex Analysis\par}
\vspace{7mm}
{\LARGE Local calculus, coherent Hahn functions,\\[3pt] and residue theory\par}
\vspace{12mm}
{\large A proof-oriented development over $\No[i]$\par}
\vspace{6mm}
{\large September 21, 2026\par}
\vspace{12mm}
\begin{abstract}
The algebraically closed proper-class field $\K=\No[i]$ admits a rich local analytic calculus, but its fine topology does not support a literal transplantation of global complex analysis. We develop two complementary theories. First, every formal power series over $\K$ is realized by Hahn summation on a sufficiently small neighborhood. This yields differentiation, Cauchy--Riemann equations, inverse and implicit function theorems, local ramification, isolated zeros, formal Cauchy identities, and coordinate-invariant residues. Second, we introduce a globally supported algebra of holomorphic functions with Hahn coefficients and evaluate it on infinitesimal neighborhoods of ordinary complex domains. In this coherent class we prove Cauchy formulas, an identity theorem, open mapping and maximum-modulus principles, Hahn--Weierstrass preparation and division, exact lifting of zero multiplicities, Rouch\'e stability, a residue theorem with genuinely surcomplex poles, and the argument principle. Affine charts transport these results to infinitesimal and infinitely large scales.

The preparation theorem is proved by well-founded coefficient recursion, not by an unjustified completeness or compactness argument. Counterexamples explain why ordinary convergence, arbitrary fine-analytic continuation, and naive Liouville theorems fail. Established surreal-analysis foundations are credited explicitly; the presentation is a rigorous synthesis and development, not a claim of priority or a solution of every global extension problem.
\end{abstract}
\vfill
{\small\noindent Prerequisites: a first course in complex analysis, elementary field theory, and familiarity with formal power series. The normal-form and real-closedness theorems for surreal numbers and the Hahn--Neumann support lemma are the foundational inputs. All subsequent surcomplex theorems are proved.}
\end{titlepage}
\tableofcontents
\newpage

\section{The problem and the architecture of the theory}

A surcomplex number is $x+iy$, where $x,y\in\No$ and $i^2=-1$. The algebra is straightforward once the real-closedness of $\No$ is known. Analysis is not. An analytic theory must specify at least three things: its notion of infinite sum, its admissible functions, and its meaning of integration. Different choices lead to different valid theorems.

The distinction driving this article is
\[
\boxed{\text{local Hahn analyticity}\quad\ne\quad\text{global analytic coherence}.}
\]
Local Hahn analyticity is deliberately broad. In fact, every formal series becomes meaningful at sufficiently small surcomplex arguments. Global coherence is an additional condition that ties different infinitesimal neighborhoods to a single ordinary complex domain through holomorphic coefficient functions. It is this condition, not the fine topology alone, that supports useful global theorems.

\subsection{Established foundations and the scope of this article}
Conway and Gonshor provide the algebraic and normal-form foundations of surreal numbers \cite{Conway,Gonshor}. Alling's monograph already develops analysis over surreal number fields, including analytic functions \cite{Alling}. Berarducci and Mantova formulate surreal analyticity using summable power series and prove their local realization and compatibility with composition; they also develop a much deeper theory of transseries as surreal functions \cite{BMgerms}. Our first part is a complexified development of this established local framework, with a separate proof using separated Hahn supports.

Function differentiation must be distinguished from a derivation on the field of surreal numbers. The latter is the subject of \cite{BMderiv}; in the present article a surcomplex coefficient is constant with respect to the analytic variable, even when it contains $\omega$. Work on alternative surreal topologies and integration is represented by \cite{RSS,CostinEhrlich}. Hahn-meromorphic functions in the sense of M\"uller and Strohmaier provide another related, but different, generalized-series theory \cite{MS}: their analytic variable is not the full surcomplex field used here.

For current context, Mantova's 2026 preprint proves monotonicity and a Taylor approximation theorem for omega-series \cite{Mantova2026}. We do not describe these as unresolved questions. None of that preprint's new results is needed below.

The specific global algebra used here is defined in Section~\ref{sec:coherent}. Its preparation, zero-lifting, and residue results are proved from explicit support calculations. These statements are the main development of the article. No assertion that their precise formulations have never appeared elsewhere is made.

\subsection{A guide to the principal results}
The local theory identifies analytic germs at a point with $\K[[X]]$. A nonconstant germ has the exact normal form
\[
 f(a+h)=f(a)+\psi(h)^m,
\]
where $\psi$ is locally invertible. Thus finite-order local branching survives essentially unchanged.

The global theory starts with expressions
\[
 F(z)=\sum_{\gamma\in S}t^\gamma f_\gamma(z),
 \qquad t^\gamma:=\omega^{-\gamma},
\]
where $S\subseteq\No$ is a well-ordered set and every $f_\gamma$ is holomorphic on the same ordinary domain. At a finite surcomplex point $a+\epsilon$, the evaluation is
\[
 \ev F(a+\epsilon)
 =\sum_{\gamma\in S}\sum_{n\geq0}
 t^\gamma\frac{f_\gamma^{(n)}(a)}{n!}\epsilon^n.
\]
The double family, not merely each inner sum, is summable.

If $F$ has leading coefficient $f_0$ after normalization, then every ordinary zero of $f_0$ lifts to exactly as many surcomplex zeros as its multiplicity. They can split, move through several infinitesimal scales, and have infinite residues. Nevertheless, the residue theorem and argument principle remain exact for the coefficientwise contour functional defined below.

\section{Surcomplex algebra, normal forms, and summation}\label{sec:foundations}

\subsection{Field structure, modulus, and valuation}
We work in a class theory such as NBG. All supports, coefficient sequences, and indexing families appearing in a single sum are sets. The word ``field'' for $\K$ is shorthand for a proper class with field operations. No sum indexed by all ordinals is used.

The standard normal-form theorem says that every surcomplex number has a unique expression
\begin{equation}\label{eq:normal}
 z=\sum_{\gamma\in S}c_\gamma t^\gamma,
 \qquad c_\gamma\in\C\setminus\{0\},
 \quad t^\gamma=\omega^{-\gamma},
\end{equation}
with $S\subseteq\No$ a well-ordered set. This follows by combining the real normal forms of the real and imaginary parts. The notation $t^\gamma$ refers to Conway monomials: $t^{\alpha+\beta}=t^\alpha t^\beta$.

For $z\ne0$, define
\[
 \vH(z)=\min\supp z,\qquad
 \lc(z)=c_{\vH(z)},\qquad \vH(0)=+\infty.
\]
Thus larger valuation means smaller magnitude. We have
\[
 \vH(zw)=\vH(z)+\vH(w),\qquad
 \vH(z+w)\geq\min\{\vH(z),\vH(w)\}.
\]
Put $\bar z=x-iy$ and $|z|=\sqrt{x^2+y^2}\in\No_{\geq0}$. Real-closedness supplies the nonnegative square root. The usual algebraic proofs give
\[
 |zw|=|z||w|,\quad |z+w|\leq |z|+|w|,\quad
 \max\{|x|,|y|\}\leq|x+iy|\leq |x|+|y|.
\]
For example, the triangle inequality follows by squaring and using
$\operatorname{Re}(z\bar w)\leq|z||w|$; this last inequality is equivalent to the elementary two-dimensional Cauchy--Schwarz inequality over a real-closed field.

Because $\No$ is real closed, adjoining $i$ gives an algebraically closed field \cite{Conway,Gonshor}. Accordingly every nonconstant polynomial over $\K$ splits, and every nonzero $a\in\K$ has an $m$th root. This is an algebraic input, not an analytic consequence of Liouville's theorem.

\subsection{Finite numbers and infinitesimals}
Define
\[
 \fin=\{z\in\K:|z|\leq n\text{ for some }n\in\N\},\qquad
 \infinit=\{z\in\K:|z|<1/n\text{ for every }n\geq1\}.
\]
Equivalently, $\fin=\{\vH\geq0\}$ and $\infinit=\{\vH>0\}$, with zero included. Every $z\in\fin$ has a unique decomposition
\begin{equation}\label{eq:st}
 z=a+\epsilon,\qquad a\in\C,\quad\epsilon\in\infinit.
\end{equation}
The map $\st(z)=a$ extracts the coefficient of $t^0$ and is a ring homomorphism $\fin\to\C$ with kernel $\infinit$.

We write $u\prec v$ when $v\ne0$ and $|u/v|$ is infinitesimal. For a subfield $k\subseteq\K$ containing $\C$, the expression $h\prec_k1$ means $|h|<r$ for every positive $r\in k\cap\No$.

\subsection{The exact meaning of a Hahn sum}
\begin{definition}[Summability]
A set-indexed family $(z_j)_{j\in J}$ is Hahn-summable when its combined support
$\bigcup_j\supp(z_j)$ is well ordered and each exponent belongs to the support of only finitely many $z_j$. Its sum is obtained by adding coefficients at each exponent.
\end{definition}
The finiteness requirement is essential. An ordinary convergent numerical series such as $\sum_{n\geq0}2^{-n}$ is not a Hahn-summable family of constants: infinitely many terms occur at exponent zero. Conversely, $\sum_{n\geq0}n!t^n$ is Hahn-summable despite the factorial coefficients.

We use the following standard support theorem \cite{Neumann}; it is the one combinatorial summation result taken as foundational.
\begin{lemma}[Hahn--Neumann support lemma]\label{lem:neumann}
Let $G$ be an ordered abelian group. If $S,T\subseteq G$ are well ordered, then $S+T$ is well ordered and each element has only finitely many representations $s+t$ with $(s,t)\in S\times T$. If $A\subseteq G_{>0}$ is well ordered, its additive monoid
\[
 \langle A\rangle=\{a_1+\cdots+a_n:n\geq0,\ a_j\in A\}
\]
is well ordered, and a given element has only finitely many finite multisets of elements of $A$ summing to it.
\end{lemma}
Ordering the members of each multiset also gives finiteness of the corresponding words, for each fixed total exponent. The strict positivity of $A$ cannot be dropped.

It follows that products of summable families are summable, with the expected product of sums. If the entire double family is summable, one may interchange its two summations. We will justify each substantial use by locating its support in $S+\langle A\rangle$, with $A$ positive, or by a separated-support argument.

\begin{proposition}[Formal series at infinitesimals]\label{prop:complexcoeff}
For $\epsilon\in\infinit$ and arbitrary complex numbers $c_n$, the family $(c_n\epsilon^n)_{n\geq0}$ is summable. Formal addition, multiplication, and substitution of zero-constant-term power series agree with evaluation.
\end{proposition}
\begin{proof}
Take $A=\supp\epsilon\subseteq\No_{>0}$. Expanding the powers gives support in $\langle A\rangle$. The finite-representation conclusion of Lemma~\ref{lem:neumann} proves summability of the fully expanded family, including its degree index. The asserted identities are then rearrangements of that family. No bound on $|c_n|$ is relevant to a Hahn support.
\end{proof}

\section{Why the ordinary topological template fails}\label{sec:topology}
The \emph{fine topology} on $\K$ has balls
$B(a,r)=\{z:|z-a|<r\}$ for every positive surreal radius $r$. Limits and derivatives in the local theory use all such radii.

A basic property of $\No$ is that between two separated sets of numbers there is another surreal number. In particular, every set of positive surreal numbers has a strictly positive lower bound. This produces behavior quite unlike $\R$ or $\C$.

\begin{proposition}[Set-sized subsets and nets]\label{prop:setdiscrete}
Every set-sized subset of $\K$ is closed and discrete in the fine topology. A net with a set-sized index set can converge in this topology only if it is eventually equal to its limit. Consequently every continuous map from the ordinary real interval $[0,1]$ into the fine surcomplex space is constant.
\end{proposition}
\begin{proof}
For a set $E\subseteq\K$ and a point $a$, the nonzero distances
$D=\{|a-e|:e\in E,\ e\ne a\}$ form a set of positive surreals. Choose $0<\delta<d$ for all $d\in D$. Then $B(a,\delta)$ meets $E$ only possibly at $a$. This proves both closedness and discreteness.

For a convergent set-indexed net, apply the same argument to the set of its values together with its proposed limit. The resulting ball forces eventual equality. Finally, the image of a map from $[0,1]$ is a set, hence has discrete subspace topology. A continuous image of a connected interval in a discrete space has at most one point.
\end{proof}
The obstruction to nonconstant set-length convergent sequences is also discussed in the surreal-analysis literature \cite{RSS}. Our use of Hahn summation explicitly avoids identifying a sum with a fine-topological limit of ordinary partial sums.

\begin{proposition}[Clopen separation]\label{prop:clopen}
For each $a\in\K$ and $r>0$, the class
\[
 M(a,r)=\{z:|z-a|\prec r\}
\]
is open and closed. Such classes separate any two distinct points.
\end{proposition}
\begin{proof}
Inside $M(a,r)$, adding any increment infinitesimal relative to $r$ stays inside. If $z\notin M(a,r)$, then $|z-a|\geq r/n$ for some positive integer $n$. Every $w$ with $|w-z|<r/(2n)$ remains outside $M(a,r)$. Thus the complement is open. Given $a\ne b$, take $r=|b-a|$.
\end{proof}
In particular, the fine topology is totally separated. Replacing ``connected complex domain'' by ``connected fine-open surcomplex domain'' would make many classical theorems vacuous.

\begin{example}[Bounded, everywhere locally constant, nonconstant]\label{ex:badliouville}
Define on all of $\K$
\[
 L(z)=\begin{cases}0,&z\in\fin,\\1,&z\notin\fin.\end{cases}
\]
A radius-one neighborhood never changes finiteness, so $L$ is locally constant. It is therefore locally represented by a constant power series, has derivative zero everywhere, and satisfies $|L|\leq1$, yet is not constant. Thus neither an unrestricted Liouville theorem nor ``derivative zero implies globally constant'' follows from local power-series analyticity.
\end{example}

\begin{example}[Agreement on $\C$ is not uniqueness]
If $f$ is ordinarily holomorphic on $U\subseteq\C$, then the function
$z\mapsto f(\st z)$ on $\halo U=\st^{-1}(U)$ is locally constant. It agrees with $f$ on $U$. Its derivative is zero, unlike the Taylor extension developed later. Agreement on the ordinary points alone does not characterize an arbitrary fine-analytic extension.
\end{example}
These examples are not reasons to abandon analysis. They identify exactly what needs to be added: summation rules locally, and coherence across ordinary base points globally.

\section{Every formal surcomplex series has a local realization}\label{sec:realization}

\subsection{Separating the variable from its coefficients}
Let $\Gamma\subseteq(\No,+)$ be a set-sized subgroup. Its Hahn field
\[
 \smallfield=\C((t^\Gamma))\subseteq\K
\]
consists of all series whose supports lie in $\Gamma$. A set of surcomplex coefficients is contained in some such field: take the subgroup generated by the union of their supports.

Since $\Gamma$ is a set, choose $\beta\in\No$ larger than every element of $\Gamma$. For $0<|h|<t^\beta$, its leading exponent $\eta=\vH(h)$ satisfies $\eta\geq\beta>\Gamma$. Thus the variable is smaller than every nonzero scale in its coefficient field.

\begin{lemma}[Separated evaluation]\label{lem:separated}
Let $P(X)=\sum_{n\geq0}a_nX^n\in\smallfield[[X]]$ and suppose $\eta=\vH(h)>\Gamma$. Then $(a_nh^n)_n$ is summable. If $m$ is the first index with $a_m\ne0$, then
\begin{equation}\label{eq:leadingevaluation}
 P(h)=a_mh^m(1+\theta),\qquad\theta\prec_{\smallfield}1.
\end{equation}
Moreover, any finite collection of formal algebraic identities involving sums, products, and compositions with zero constant term is respected by evaluation at sufficiently small $h$.
\end{lemma}
\begin{proof}
Write $h=c\,t^\eta(1+u)$, with $c\in\C^\times$ and $u\in\infinit$. Before introducing $u$, the combined support is
\[
 B=\bigcup_{n\geq0}\bigl(n\eta+\supp a_n\bigr).
\]
Every exponent in the $n$th block is smaller than every exponent in the $(n+1)$st block: their difference has the form $\eta+\gamma$ with $\gamma\in\Gamma$, and is positive. Hence $B$ is an ordered sum of well-ordered blocks and is well ordered. A member of $B$ determines its block uniquely.

Expand $(1+u)^n$ by the finite binomial formula. The full support is contained in $B+\langle\supp u\rangle$. Lemma~\ref{lem:neumann} gives finite decompositions into a base exponent and a word of positive exponents. The base exponent fixes $n$, so only finitely many binomial terms contribute. This proves summability of the fully expanded family.

For $n>m$, the leading exponent of $(a_n/a_m)h^{n-m}$ exceeds every member of $\Gamma$. The first nonzero term of their sum therefore does also, giving \eqref{eq:leadingevaluation}. Formal products follow by the product rule for summable families. For compositions, the inner series has positive formal order, so only finitely many coefficient expressions contribute at a fixed formal degree. All resulting coefficients remain in $\smallfield$, and the same separated-support proof justifies regrouping by total degree.
\end{proof}

\begin{theorem}[Universal local realization]\label{thm:realization}
Every $P\in\K[[X]]$ defines a function by Hahn evaluation on a ball $B(0,r)$ with $r>0$. A nonzero series gives a nonzero germ. Evaluation identifies $\K[[X]]$ with the algebra of germs at zero defined by local Hahn power series; the identification respects differentiation and composition when the inner germ vanishes at zero.
\end{theorem}
\begin{proof}
Put the coefficient sequence inside a set-sized $\smallfield$ and choose $r=t^\beta$ as above. Lemma~\ref{lem:separated} supplies evaluation. Its leading-term assertion proves injectivity: a nonzero series cannot vanish at all sufficiently small nonzero arguments. Surjectivity is the definition of the germ algebra. Differentiation is proved in the next section, independently of any sequential convergence claim.
\end{proof}

\subsection{Several variables and recentering}
The same result holds for formal series in any fixed finite number of variables. Here is the support argument needed later. If all nonzero $h_j$ have valuation greater than $\Gamma$, choose $\eta=\min_j\vH(h_j)$ and write $h_j=t^\eta v_j$ with each $v_j$ finite. Separate the constant terms of the finitely many $v_j$ from their positive-support tails. Grouping a formal series by total degree gives finitely many monomials at each degree, the separated blocks $n\eta+\Gamma$, and an additional monoid generated by finitely many well-ordered positive supports. Lemma~\ref{lem:neumann} again proves summability of the full expansion.

In particular, if $x$ and $h$ are sufficiently small relative to a coefficient field, the binomial expansion and regrouping give
\begin{equation}\label{eq:recenter}
 P(x+h)=\sum_{k\geq0}\frac{P^{(k)}(x)}{k!}h^k,
 \qquad
 \frac{P^{(k)}(x)}{k!}
 =\sum_{n\geq k}\binom nk a_nx^{n-k}.
\end{equation}
The full double family is controlled by the preceding two-variable argument. Consequently a realized series is locally a power series at every point of a sufficiently small realization ball, not only at its original center.

\begin{example}[No complex radius is required]\label{ex:euler}
The series $E(X)=\sum_{n\geq0}n!X^n$ has zero ordinary complex radius of convergence, but defines a surcomplex function at every $\epsilon\in\infinit$. Its formal identity
\[
 X^2E'(X)+(X-1)E(X)=-1
\]
therefore becomes an exact functional differential equation there. This does not identify $E$ with any particular Borel sum on a classical sector.
\end{example}

\begin{example}[An essentially surreal convergence scale]
The series $P(X)=\sum_{n\geq0}\omega^{n^2}X^n$ is summable at $X=\omega^{-\omega}$: its exponents in $t$ are $n\omega-n^2$, a strictly increasing sequence. It is not summable at $X=\omega^{-1}$, where the exponents $n-n^2$ contain an infinite decreasing sequence. The local scale can depend on the surreal sizes of the coefficients, not merely on their ordinary numerical growth.
\end{example}

\section{Differentiation and Cauchy--Riemann equations}\label{sec:derivatives}

\begin{definition}[Hahn-analytic functions]
A function $f:V\to\K$ on a fine-open class $V\subseteq\K$ is \emph{Hahn analytic} if, near every $a\in V$, it is the Hahn evaluation of a formal power series in $z-a$. Its derivative, when it exists, is the fine limit
\[
 f'(a)=\lim_{h\to0}\frac{f(a+h)-f(a)}h.
\]
The limit quantifies over every positive surreal tolerance and sufficiently small nonzero surcomplex $h$.
\end{definition}

\begin{theorem}[Differentiation, Taylor coefficients, and calculus]\label{thm:derivative}
A Hahn-analytic function has derivatives of every finite order. If
$f(a+h)=\sum_{n\geq0}a_nh^n$ near $a$, then
\[
 f^{(k)}(a)=k!a_k,
 \qquad
 f'(a+h)=\sum_{n\geq1}na_nh^{n-1}
\]
on a sufficiently small neighborhood. The product, quotient, and chain rules hold wherever the indicated functions are defined and the denominator is nonzero.
\end{theorem}
\begin{proof}
At the center, write
\[
 P(h)=a_0+a_1h+h^2R(h),\qquad R(X)=a_2+a_3X+\cdots.
\]
Place all coefficients in a set-sized Hahn field. For $h$ sufficiently small relative to that field, Lemma~\ref{lem:separated} gives $|R(h)|\leq M:=|a_2|+1$. Given any surreal tolerance $\delta>0$, impose also $|h|<\delta/M$. Then
\[
 \left|\frac{P(h)-P(0)}h-a_1\right|
 =|hR(h)|<\delta.
\]
Thus $P'(0)=a_1$. At another sufficiently small point, use the recentering identity \eqref{eq:recenter}. Repeating the argument proves all derivative formulas.

The product rule follows either from the finite difference-quotient identity or from formal multiplication. A series with nonzero constant term has a formal reciprocal, obtained from the geometric series after factoring its constant term; realization proves the quotient rule. For composition, first translate the inner function's value to zero, use formal composition, and differentiate the resulting formal identity. Continuity needed in the limit arguments also follows from the same linear-plus-bounded-remainder estimate.
\end{proof}

\begin{remark}[A derivative of functions, not of constants]
If $f(z)=\omega z^2$, then $f'(z)=2\omega z$. The coefficient $\omega$ is constant in this derivative. A surreal-field derivation for which $\partial\omega=1$ is a different operation. Conflating these operations would invalidate even elementary calculations.
\end{remark}

\subsection{The differential characterization}
Identify $\K$ with $\No^2$. A map $f=u+iv$ is \emph{totally differentiable over $\No$} at $(x,y)$ if
\[
 f((x,y)+(h,k))=f(x,y)+L(h,k)+o\bigl(\sqrt{h^2+k^2}\bigr)
\]
for an $\No$-linear map $L:\No^2\to\No^2$.

\begin{proposition}[Cauchy--Riemann criterion]\label{prop:CR}
For a totally differentiable map $f=u+iv$, surcomplex differentiability at a point is equivalent to
\[
 u_x=v_y,\qquad u_y=-v_x.
\]
In that case $f'=u_x+iv_x$. Every Hahn-analytic function satisfies these equations. Its real and imaginary parts are harmonic:
\[
 u_{xx}+u_{yy}=0,\qquad v_{xx}+v_{yy}=0.
\]
\end{proposition}
\begin{proof}
A surcomplex-linear map is multiplication by $A+iB$, whose real matrix is
$\left(\begin{smallmatrix}A&-B\\B&A\end{smallmatrix}\right)$. The displayed equations say precisely that the total derivative has this form. The total remainder estimate is then exactly the complex difference-quotient estimate. Hahn power-series expansions give total differentiability of all orders. Differentiate the Cauchy--Riemann equations once more and commute formal mixed derivatives to obtain the Laplace equations.
\end{proof}
Existence of partial derivatives alone is not the hypothesis: total differentiability is required, just as in a careful classical statement.

\begin{theorem}[Cauchy--Riemann equations for real-analytic germs]\label{thm:analyticCR}
Suppose $u$ and $v$ have local Hahn-summable formal expansions in the two real variables $x$ and $y$. Then $u+iv$ is Hahn analytic in $z=x+iy$ if and only if its Cauchy--Riemann equations hold as local identities.
\end{theorem}
\begin{proof}
Translate the point to the origin. In $\K[[X,Y]]$, make the invertible linear change
$Z=X+iY$, $W=X-iY$. The Cauchy--Riemann equations become
$\partial F/\partial W=0$. In characteristic zero this means that every coefficient containing a positive power of $W$ vanishes. Hence $F\in\K[[Z]]$.

To pass from an identity of evaluated real-analytic germs to an identity of formal series, note that formal evaluation in finitely many real variables is faithful. Indeed, take the least nonzero homogeneous part of a nonzero formal series and choose a rational direction on which that polynomial is nonzero. Such a direction exists because a nonzero polynomial over a characteristic-zero field cannot vanish on all rational tuples. Restriction to that line gives a nonzero one-variable series, contradicting Lemma~\ref{lem:separated} if the evaluated germ were zero. The converse follows from Proposition~\ref{prop:CR}.
\end{proof}

\section{Inverse functions, implicit functions, and local geometry}\label{sec:localgeometry}

\begin{theorem}[Inverse function theorem]\label{thm:inverse}
If $f$ is Hahn analytic near $a$ and $f'(a)\ne0$, then there are fine-open neighborhoods $V$ of $a$ and $W$ of $f(a)$ on which $f:V\to W$ is bijective with a Hahn-analytic inverse.
\end{theorem}
\begin{proof}
Translate both centers to zero and write $F(X)=a_1X+a_2X^2+\cdots$ with $a_1\ne0$. Construct $G(Y)=\sum_{n\geq1}b_nY^n$ recursively by
\[
 b_1=a_1^{-1},\qquad
 b_n=-a_1^{-1}[Y^n]\sum_{k=2}^{n}a_k
       \left(\sum_{j=1}^{n-1}b_jY^j\right)^k\quad(n\geq2).
\]
This gives $F\circ G=Y$. The same coefficient argument gives uniqueness, and formal composition then gives $G\circ F=X$ (or construct the left inverse and compare). All coefficients constitute a set, so both series and both identities are realized on sufficiently small neighborhoods.

Choose initial balls where these identities hold. By continuity shrink a target ball $W$ so that $G(W)$ lies in the source ball. Intersect the source ball with $F^{-1}(W)$. This is open, and the two identities give the desired inverse maps on these neighborhoods.
\end{proof}

\begin{theorem}[Implicit function theorem]\label{thm:implicit}
Let $F$ be a vector of $q$ Hahn-analytic functions of $p+q$ surcomplex variables near $(a,b)$, with $F(a,b)=0$. If the $q\times q$ Jacobian $\partial F/\partial y$ is invertible there, then, near $(a,b)$, the solutions of $F(x,y)=0$ are exactly $y=g(x)$ for a unique Hahn-analytic germ $g$ with $g(a)=b$.
\end{theorem}
\begin{proof}
For a formal vector map with invertible linear part, construct its inverse degree by degree. At homogeneous degree $n$, all lower-degree terms are known, and the unknown vector of degree-$n$ coefficients is obtained by multiplying a known vector by the inverse linear matrix. This proves the formal inverse theorem in finitely many variables. Apply it to $(x,y)\mapsto(x,F(x,y))$. Its inverse at $(x,0)$ supplies $g$. The several-variable evaluation theorem realizes the formal identities. They give both existence and uniqueness after shrinking neighborhoods as in Theorem~\ref{thm:inverse}.
\end{proof}

\begin{theorem}[Local ramification normal form]\label{thm:normalform}
If the germ of $f-f(a)$ is nonzero, there are a unique integer $m\geq1$ and an invertible Hahn-analytic germ $\psi$ with $\psi(0)=0$ such that
\begin{equation}\label{eq:ramification}
 f(a+h)=f(a)+\psi(h)^m.
\end{equation}
The integer $m$ is the order of the first nonzero Taylor coefficient of $f-f(a)$.
\end{theorem}
\begin{proof}
Factor the formal expansion as
$F(X)-F(0)=a_mX^m(1+R(X))$, with $R(0)=0$ and $a_m\ne0$. Choose $c\in\K$ with $c^m=a_m$ and set
\[
 \psi(X)=cX\sum_{n\geq0}\binom{1/m}{n}R(X)^n.
\]
The formal binomial identity gives \eqref{eq:ramification}, and $\psi'(0)=c\ne0$. Apply realization and the inverse theorem. Uniqueness of $m$ follows from the least nonzero formal degree; $\psi$ itself is not unique because of the choice of an $m$th root.
\end{proof}

\begin{corollary}[Isolated zeros and local open mapping]\label{cor:localmapping}
A Hahn-analytic germ which is not identically zero has an isolated zero at its center, if it has a zero there at all. At any point where the germ of $f$ is nonconstant, $f$ maps every sufficiently small neighborhood onto a set containing a neighborhood of its value. Hence a Hahn-analytic map with nonconstant germ at every point is open.
\end{corollary}
\begin{proof}
For isolated zeros factor $f(a+h)=h^mU(h)$ with $U(0)\ne0$; $U$ is nonzero near zero. For local openness use \eqref{eq:ramification}. The map $w\mapsto w^m$ takes $B(0,r)$ onto $B(0,r^m)$: roots exist in $\K$, and their moduli satisfy $|w|^m=|w^m|$. The coordinate $\psi$ is locally invertible.
\end{proof}

\begin{corollary}[Local maximum-modulus principle]\label{cor:localmaximum}
If $|f|$ has a fine-local maximum at $a$, then the germ of $f$ at $a$ is constant.
\end{corollary}
\begin{proof}
Otherwise use \eqref{eq:ramification}. If $f(a)=0$, nearby nonzero values contradict a maximum. If $b=f(a)\ne0$, choose an arbitrarily small positive surreal $s$ and an $m$th root $w$ of $sb$ sufficiently small to lie in the inverse-coordinate neighborhood. At $h=\psi^{-1}(w)$ one has $f(a+h)=b(1+s)$, whose modulus is larger than $|b|$.
\end{proof}
These statements concern \emph{germs}. A globally nonconstant locally constant function, such as Example~\ref{ex:badliouville}, is not an exception to them: it has a constant germ at each point.

\begin{example}[An inverse with infinite coefficients]
For $f(z)=z+\omega z^2$ near zero,
\[
 f^{-1}(w)=\frac{\sqrt{1+4\omega w}-1}{2\omega}
 =\sum_{n\geq1}(-1)^{n-1}C_{n-1}\omega^{n-1}w^n,
\]
where $C_j=\frac1{j+1}\binom{2j}{j}$. The square root is the local branch with constant term one. The identity is valid whenever $\omega w$ is infinitesimal, and in particular on a fine neighborhood of zero.
\end{example}

\section{Local meromorphic functions and formal residues}\label{sec:formalresidues}

\subsection{Meromorphic germs and primitives}
The field of fractions of $\K[[X]]$ is the ordinary formal Laurent-series field $\K((X))$: every nonzero series is $X^mU$ with $U$ a unit. Its members have only finitely many negative powers. They evaluate on sufficiently small punctured neighborhoods.

\begin{definition}[Local residue]
If a meromorphic germ at $a$ has expansion
$f(a+X)=\sum_{n\geq-N}c_nX^n$, define
\[
 \Res_a(f(z)\dd z)=c_{-1}.
\]
The residue is a surcomplex number. Its size need not be finite.
\end{definition}

\begin{theorem}[Primitives and logarithmic residues]\label{thm:formalprimitive}
A meromorphic germ has a meromorphic primitive if and only if its residue is zero. If $f\ne0$ is meromorphic, then
\[
 \Res_a\left(\frac{f'}f\dd z\right)=\ord_a f\in\Z.
\]
In particular, every analytic germ has an analytic primitive.
\end{theorem}
\begin{proof}
Differentiation of a Laurent series never produces an $X^{-1}$ term. Conversely, when that term is absent, integrate coefficientwise:
$\sum_{n\ne-1}c_nX^{n+1}/(n+1)$. This is again a Laurent series. If $f=X^mU$ with $U$ a unit, then $f'/f=m/X+U'/U$, and $U'/U$ is a power series.
\end{proof}
Thus the only local obstruction to a Laurent primitive is exactly the classical one. A formal logarithm can be adjoined to integrate $c_{-1}/X$, but a global single-valued logarithm is an additional problem.

\begin{theorem}[Residues under a change of coordinate]\label{thm:reschange}
Let $\phi(X)\in X\K[[X]]$ have order $e\geq1$, and let $F\in\K((Y))$. Then
\[
 \Res_{X=0}F(\phi(X))\phi'(X)\dd X
 =e\,\Res_{Y=0}F(Y)\dd Y.
\]
For $e=1$, the residue of a meromorphic differential is invariant under local Hahn-analytic coordinate changes.
\end{theorem}
\begin{proof}
For the monomial $Y^n$, $n\ne-1$, the substituted differential is the derivative of $\phi^{n+1}/(n+1)$ and has residue zero. For $Y^{-1}$, write $\phi=X^eU$; then $\phi'/\phi=e/X+U'/U$ has residue $e$. Linearity completes the proof. There is no problematic exchange of infinite sums: nonnegative powers of $Y$ cannot contribute a negative power after substitution and multiplication by $\phi'$, and only finitely many negative powers occur in $F$.
\end{proof}

\subsection{A Cauchy formula before contours}
There is also an exact formal Cauchy identity. For $F(X)\in\K[[X]]$, expand the kernel in the indicated order:
\[
 \frac1{X-Y}=\sum_{n\geq0}Y^nX^{-n-1}
 \quad\text{in }\K((X))[[Y]].
\]
Then
\begin{equation}\label{eq:formalCauchy}
 \Res_{X=0}\frac{F(X)}{X-Y}\dd X=F(Y),
 \qquad
 \Res_{X=0}\frac{F(X)}{(X-Y)^{k+1}}\dd X
 =\frac{F^{(k)}(Y)}{k!}.
\end{equation}
Indeed, take the residue coefficientwise in $Y$. The $n$th coefficient in the first identity is $[X^n]F$. The second follows by differentiating the first $k$ times in $Y$.

Formula~\eqref{eq:formalCauchy} is a coefficient-extraction theorem, not an integral over an unspecified surreal circle. It holds for every local Hahn-analytic germ, including the factorial series of Example~\ref{ex:euler}. A contour interpretation for a more structured global class is established in Section~\ref{sec:contours}.

\section{Canonical lifts of ordinary holomorphic functions}\label{sec:canonical}

For an ordinary open set $U\subseteq\C$, define its \emph{finite halo}
\[
 \halo U=\{a+\epsilon:a\in U,\ \epsilon\in\infinit\}=\st^{-1}(U).
\]
This is fine-open, but it is not connected in the fine topology even when $U$ is a disk. The word ``connected'' applied to $U$ below always means ordinary complex connectedness.

\begin{definition}[Canonical Taylor lift]
For $f\in\hol(U)$ set
\begin{equation}\label{eq:canonicallift}
 \ev f(a+\epsilon)=\sum_{n\geq0}\frac{f^{(n)}(a)}{n!}\epsilon^n.
\end{equation}
The decomposition \eqref{eq:st} is unique, and Proposition~\ref{prop:complexcoeff} makes the sum well defined.
\end{definition}

\begin{theorem}[Functoriality of the lift]\label{thm:canonicalfunctor}
Canonical lifting preserves addition, multiplication, constants, and the identity function. It commutes with differentiation. If $g:U\to V$ and $f:V\to\C$ are holomorphic, then
\[
 \ev{f\circ g}=\ev f\circ\ev g.
\]
Reciprocals are preserved where the original function is nonzero. Moreover $\st(\ev f(z))=f(\st z)$.
\end{theorem}
\begin{proof}
At a fixed $a$, the claims about sums, products, derivatives, and reciprocals are the corresponding identities of ordinary Taylor series, evaluated at an infinitesimal. Write $\ev g(a+\epsilon)=g(a)+\delta$; all nonconstant Taylor terms are infinitesimal, so $\delta\in\infinit$. The ordinary formal chain identity and Proposition~\ref{prop:complexcoeff} give the composition assertion. Extracting exponent zero gives the standard-part assertion.
\end{proof}

\begin{example}[Elementary functions]
On the appropriate ordinary domains, the lifts satisfy
\[
 \ev{\exp}(a+\epsilon)=e^a\sum_{n\geq0}\frac{\epsilon^n}{n!},\qquad
 \ev{\log}(a+\epsilon)=\log a+
 \sum_{n\geq1}\frac{(-1)^{n+1}}n\left(\frac\epsilon a\right)^n.
\]
For the logarithm, $a$ lies in the chosen ordinary branch domain and is nonzero. The lifted trigonometric functions satisfy the usual algebraic addition identities. These are functions on finite halos, not a declaration of a canonical exponential or trigonometric parametrization on all of $\K$.
\end{example}

\subsection{A Schwarz lemma and a Schwarz--Pick theorem}
Let $\mathbb D=\{a\in\C:|a|<1\}$. Notice that $\halo{\mathbb D}$ excludes points with $|z|<1$ but $\st|z|=1$. It is the halo of the \emph{open standard disk}, not the entire fine-open surreal unit disk.

\begin{theorem}[Lifted Schwarz lemma]\label{thm:Schwarz}
If $f:\mathbb D\to\mathbb D$ is holomorphic and $f(0)=0$, then
\[
 |\ev f(z)|\leq |z|\qquad(z\in\halo{\mathbb D}).
\]
Unless $f(z)=cz$ with $|c|=1$, the inequality is strict for every nonzero $z$ in that halo.
\end{theorem}
\begin{proof}
For $a=\st z\ne0$, classical Schwarz gives $|f(a)|<|a|$ unless $f$ is a rotation. A strict real gap survives infinitesimal perturbation. If $a=0$, then $z$ is infinitesimal and
$\ev f(z)=z(f'(0)+\eta)$ with $\eta$ infinitesimal. For a nonrotation $|f'(0)|<1$, again yielding strict inequality. For a rotation, the lifted identity is exact. The classical Schwarz lemma used here is the usual consequence of the maximum principle applied to $f(z)/z$.
\end{proof}

\begin{theorem}[Lifted Schwarz--Pick inequality]\label{thm:SchwarzPick}
If $f:\mathbb D\to\mathbb D$ is holomorphic, then for $z,w\in\halo{\mathbb D}$,
\[
 \left|\frac{\ev f(z)-\ev f(w)}{1-\overline{\ev f(w)}\ev f(z)}\right|
 \leq
 \left|\frac{z-w}{1-\bar w z}\right|.
\]
For distinct $z,w$, equality occurs precisely when $f$ is a disk automorphism.
\end{theorem}
\begin{proof}
The denominators have positive nonzero real standard-part moduli. If $\st z\ne\st w$, classical Schwarz--Pick is strict for a nonautomorphism, so the strict real gap persists under lifting.

Suppose $\st z=\st w=a$ and $z\ne w$. Taylor expansion in the two infinitesimal increments gives
\[
 \frac{\ev f(z)-\ev f(w)}{z-w}=f'(a)+\eta,
 \qquad \eta\in\infinit.
\]
This factorization remains valid however small $z-w$ is: divide the formal difference of powers by $z-w$ before evaluation. The ratio of the left side of the claimed inequality to its right side has standard part
\[
 \frac{|f'(a)|(1-|a|^2)}{1-|f(a)|^2}<1
\]
for a nonautomorphism, by the differential form of classical Schwarz--Pick. Thus it is strictly less than one. If $f$ is an automorphism, its rational formula and the identity
$|1-\bar wz|^2-|z-w|^2=(1-|z|^2)(1-|w|^2)$ prove exact preservation algebraically over $\K$. The case $z=w$ is immediate.
\end{proof}
These results are restricted to canonical lifts. An arbitrary locally constant extension agreeing with $f$ on ordinary points need not satisfy the corresponding sharp infinitesimal statements.

\section{Coherent holomorphic functions with Hahn coefficients}\label{sec:coherent}

\subsection{The global algebra}
\begin{definition}[Coherent Hahn section]
For an ordinary domain $U\subseteq\C$, let $\A(U)$ consist of formal expressions
\begin{equation}\label{eq:coherent}
 F=\sum_{\gamma\in S}t^\gamma f_\gamma,
 \qquad f_\gamma\in\hol(U),
\end{equation}
whose nonzero coefficient functions have a common well-ordered set support $S\subseteq\No$. Addition and multiplication are coefficientwise addition and Hahn convolution. Define
\[
 DF=\sum_\gamma t^\gamma f_\gamma',\qquad
 \nu(F)=\min\{\gamma:f_\gamma\not\equiv0\}
\]
for $F\ne0$, and $\nu(0)=+\infty$. Write $\A_+(U)=\{F:\nu(F)>0\}$, including zero.
\end{definition}
Here ``coherent'' means the specific common-support holomorphic-coefficient condition in \eqref{eq:coherent}; it is not a claim about coherent sheaves in algebraic geometry. Restriction gives a presheaf of $\K$-algebras. We do not assert arbitrary sheaf gluing: a union of supports from an unrestricted cover need not be well ordered. All global theorems assume the displayed common-support representation on the stated domain.

Since $U$ is connected, $\hol(U)$ has no zero divisors. The leading coefficients of a product multiply, so $\A(U)$ also has no zero divisors and $\nu(FG)=\nu(F)+\nu(G)$. Constants from all of $\K$ embed in this algebra. Each individual construction takes place over a set-sized subgroup generated by its supports, even though no single subgroup is fixed for the whole theory.

\begin{definition}[Evaluation on a halo]
For $a\in U$ and $\epsilon\in\infinit$, define
\begin{equation}\label{eq:evalcoherent}
 \ev F(a+\epsilon)
 =\sum_{\gamma\in S,\ n\geq0}
 t^\gamma\frac{f_\gamma^{(n)}(a)}{n!}\epsilon^n.
\end{equation}
\end{definition}

\begin{theorem}[Coherent evaluation theorem]\label{thm:coherenteval}
The entire family in \eqref{eq:evalcoherent} is summable. Evaluation is an injective $\K$-algebra homomorphism from $\A(U)$ to functions on $\halo U$. Its values are Hahn-analytic functions, and
\[
 (\ev F)'=\ev{DF},\qquad
 \ev F(w+h)=\sum_{n\geq0}\frac{\ev{D^nF}(w)}{n!}h^n
 \quad(w\in\halo U,\ h\in\infinit).
\]
If $F=t^\beta(f_0+E)$ with $E\in\A_+(U)$ and $f_0(a)\ne0$, then
\begin{equation}\label{eq:leadingcoherent}
 \ev F(a+\epsilon)=t^\beta\bigl(f_0(a)+\eta\bigr),
 \qquad \eta\in\infinit.
\end{equation}
\end{theorem}
\begin{proof}
Let $A=\supp\epsilon\subseteq\No_{>0}$. The fully expanded support in \eqref{eq:evalcoherent} lies in $S+\langle A\rangle$. At a fixed exponent, the support lemma allows only finitely many choices of an element of $S$ and a word in $A$, and hence only finitely many degrees $n$. This proves summability without any uniform analytic bounds on the coefficient functions.

For multiplication use the same argument with the sum of two base supports. For Taylor translation write $w=a+\epsilon$ and expand $(\epsilon+h)^j$. Now the positive set is $\supp\epsilon\cup\supp h$, and all mixed terms again form a summable family. Ordinary Taylor identities at $a$ give the formula and, by Theorem~\ref{thm:derivative}, the derivative identity. At standard points evaluation is simply $\sum t^\gamma f_\gamma(a)$. If these values vanish for every $a\in U$, uniqueness of normal forms forces all coefficient functions to vanish. Thus evaluation is injective. Finally, after normalization, every term except $f_0(a)$ has strictly positive support, proving \eqref{eq:leadingcoherent}.
\end{proof}

\begin{proposition}[Inverting a nowhere-zero leading coefficient]\label{prop:unit}
If $F=t^\beta(f_0+E)$ with $E\in\A_+(U)$ and $f_0$ nowhere zero on $U$, then $F$ has an inverse in $\A(U)$, and its evaluated inverse is $1/\ev F$ throughout $\halo U$.
\end{proposition}
\begin{proof}
The inverse is
\[
 F^{-1}=t^{-\beta}f_0^{-1}\sum_{n\geq0}(-E/f_0)^n.
\]
The powers have supports in the monoid generated by the positive support of $E$. Neumann's lemma proves summability of the entire geometric family in the coefficient algebra. Multiplication gives one, and evaluation preserves this identity. Equation~\eqref{eq:leadingcoherent} also shows directly that no evaluated value of $F$ is zero.
\end{proof}

\subsection{Identity and maximum principles recovered}
\begin{theorem}[Coherent identity theorem]\label{thm:identity}
Let $U$ be connected and $F\in\A(U)$. Each of the following implies $F=0$:
\begin{enumerate}[label=(\roman*)]
\item $\ev F$ has the zero germ at one point of $\halo U$;
\item $\ev F(a)=0$ on a set of ordinary points with an ordinary accumulation point in $U$;
\item $\ev F$ has zeros whose distinct standard parts accumulate in $U$.
\end{enumerate}
Consequently two coherent sections agreeing on a fine neighborhood of a single halo point agree on the whole halo.
\end{theorem}
\begin{proof}
Suppose $F\ne0$, and let $\beta=\nu(F)$ with leading coefficient $f_\beta\not\equiv0$. At $w=a+\epsilon$, the coefficient of $t^\beta$ in $\ev{D^nF}(w)$ is exactly $f_\beta^{(n)}(a)$: positive powers of $\epsilon$ and higher Hahn indices cannot contribute at exponent $\beta$. A zero germ forces every derivative to be zero, hence all ordinary derivatives of $f_\beta$ at $a$ vanish. Its ordinary Taylor series and the classical identity theorem contradict $f_\beta\not\equiv0$.

For (ii), uniqueness of the Hahn normal form makes every coefficient function vanish on the indicated standard set, so ordinary identity applies coefficientwise. For (iii), \eqref{eq:leadingcoherent} says that the standard part of any zero must be a zero of $f_\beta$. Distinct such points accumulating in $U$ again contradict ordinary identity. Apply (i) to a difference for the last assertion.
\end{proof}

\begin{corollary}[Open mapping and maximum modulus]\label{cor:globalmapping}
If $F\in\A(U)$ is not a constant section, then $\ev F$ is an open map for the fine topologies and $|\ev F|$ has no fine-local maximum. If $DF=0$, then $F$ is a constant surcomplex section.
\end{corollary}
\begin{proof}
A constant germ at any point would make $F-c$ have zero germ for some $c\in\K$, so Theorem~\ref{thm:identity} would give $F=c$. Thus every germ is nonconstant, and Corollaries~\ref{cor:localmapping} and \ref{cor:localmaximum} apply. For the derivative statement, $DF=0$ means $f_\gamma'=0$ for each ordinary holomorphic coefficient; connectedness makes each coefficient constant.
\end{proof}
The connectedness input lies in the coefficient functions on $U$, not in the topology of $\halo U$.

\section{Coefficientwise contours and Cauchy formulas}\label{sec:contours}

\subsection{Definition, and why it is not an ordinary surreal path integral}
Let $\gamma$ be an ordinary piecewise $C^1$ complex path with compact image in $U$. For $F\in\A(U)$ define
\begin{equation}\label{eq:contourdefinition}
 \hint_\gamma F(z)\dd z
 :=\sum_{\alpha}t^\alpha\int_\gamma f_\alpha(z)\dd z.
\end{equation}
The resulting support is a subset of the support of $F$, so the sum exists. The superscript $H$ marks a \emph{coefficientwise contour functional}. A closed path is denoted by $\hoint$.

This functional is not defined as a fine-topological limit of Riemann sums in $\K$. Proposition~\ref{prop:setdiscrete} shows that a nonconstant ordinary path is not continuous when its values are viewed in the full fine surcomplex topology. Instead, the ordinary complex parameter geometry is retained in the coefficients, and Hahn summation combines their integrals. This is an explicit choice of integration structure.

Linearity over $\K$, orientation reversal, concatenation, and ordinary change of variable follow coefficientwise. For constant surcomplex scalars, the required products of summable families are justified by the support lemma. If $\gamma$ is closed and $F$ is holomorphic on a neighborhood of a simply connected region containing it, ordinary Cauchy's theorem gives
\begin{equation}\label{eq:Cauchytheorem}
 \hoint_\gamma F(z)\dd z=0.
\end{equation}
More generally the usual homological version holds whenever it holds for every coefficient function.

\begin{theorem}[Primitives on a simply connected base]\label{thm:coherentprimitive}
If $U$ is simply connected and $F\in\A(U)$, then there exists $G\in\A(U)$ with $DG=F$. Specifying $\ev G(a_0)\in\K$ at one ordinary point determines $G$ uniquely.
\end{theorem}
\begin{proof}
For each coefficient choose the ordinary holomorphic primitive vanishing at $a_0$. Their Hahn sum has support contained in that of $F$ and supplies $G$. Add the desired constant. Uniqueness follows from the last assertion of Corollary~\ref{cor:globalmapping}.
\end{proof}

\subsection{Cauchy's formula at genuinely surcomplex points}
Let $D$ be a bounded Jordan domain with positively oriented piecewise $C^1$ boundary $\gamma$, and suppose $F\in\A(U)$ for a neighborhood $U$ of $\overline D$.

For $w=a+\epsilon\in\halo D$, the kernel on a sufficiently small ordinary collar of $\gamma$ is interpreted as the coherent section
\begin{equation}\label{eq:movingkernel}
 \frac1{z-w}=\sum_{n\geq0}\frac{\epsilon^n}{(z-a)^{n+1}}.
\end{equation}
After expanding the fixed surcomplex $\epsilon$, the positive-support lemma makes this a Hahn section, and its value at ordinary contour points equals the field-theoretic reciprocal. Thus products with $F$ have a well-defined coefficientwise integral.

\begin{theorem}[Surcomplex Cauchy integral formula]\label{thm:Cauchy}
For every $w\in\halo D$ and $k\geq0$,
\begin{equation}\label{eq:Cauchy}
 \boxed{\quad
 \frac{\ev{D^kF}(w)}{k!}
 =\frac1{2\pi i}\hoint_\gamma
       \frac{F(z)}{(z-w)^{k+1}}\dd z.
 \quad}
\end{equation}
\end{theorem}
\begin{proof}
For $k=0$, multiply \eqref{eq:movingkernel} by $F$. The full family has support in $S+\langle\supp\epsilon\rangle$. Consequently the coefficientwise integral can be regrouped as
\[
 \sum_{\alpha,n}t^\alpha\epsilon^n
   \frac1{2\pi i}\int_\gamma
        \frac{f_\alpha(z)}{(z-a)^{n+1}}\dd z
 =\sum_{\alpha,n}t^\alpha\epsilon^n
        \frac{f_\alpha^{(n)}(a)}{n!}.
\]
Classical Cauchy gives the equality, and the result is \eqref{eq:evalcoherent}. For $k>0$, expand
\[
 (z-a-\epsilon)^{-k-1}
 =\sum_{n\geq0}\binom{n+k}{k}
   \epsilon^n(z-a)^{-n-k-1}.
\]
The same calculation gives
$\sum_{\alpha,n}t^\alpha\epsilon^n f_\alpha^{(n+k)}(a)/(k!n!)$,
which is the left side of \eqref{eq:Cauchy}.
\end{proof}
This formula recovers all infinitesimal Taylor data at $w$, not merely the value at its standard part.

\subsection{Coefficientwise Cauchy estimates}
Suppose the closed ordinary disk $\{z\in\C:|z-a|\leq R\}$ is contained in $U$. Set
\[
 M_\alpha(R)=\max_{|z-a|=R}|f_\alpha(z)|,
 \qquad M_F(R)=\sum_\alpha M_\alpha(R)t^\alpha.
\]
The maximum is an ordinary real maximum for each coefficient separately. No maximum of a surcomplex-valued function on a fine circle is assumed.

\begin{proposition}[Cauchy majorant]\label{prop:Cauchyestimate}
For $n\geq0$,
\[
 \left|\frac{\ev{D^nF}(a)}{n!}\right|
 \leq \frac{M_F(R)}{R^n}.
\]
\end{proposition}
\begin{proof}
Ordinary Cauchy estimates give
$|f_\alpha^{(n)}(a)|/n!\leq M_\alpha(R)R^{-n}$.
For a Hahn number $z=\sum c_\alpha t^\alpha$, let
$C=\sum |c_\alpha|t^\alpha$. Then $C\geq|z|$. Indeed,
\[
 C^2-z\bar z
 =\sum_{\alpha,\beta}
 \bigl(|c_\alpha||c_\beta|-\operatorname{Re}(c_\alpha\bar c_\beta)\bigr)
 t^{\alpha+\beta}
\]
has nonnegative real coefficients; convolution is finite at each exponent. Since $C\geq0$, the claimed comparison follows by squaring. Apply this majorant to the derivative coefficients, then use their coefficientwise inequalities.
\end{proof}
An inequality between Hahn numbers does not generally give inequalities at every exponent. The direction of the preceding argument matters: coefficientwise estimates imply an ordered-field bound, not conversely.

\subsection{A Morera criterion}
\begin{proposition}[Coherent Morera theorem]\label{prop:Morera}
Let $F=\sum t^\alpha f_\alpha$ have well-ordered set support and continuous complex coefficient functions on $U$. Suppose its coefficientwise integral around every triangle whose closed interior lies in $U$ is zero. Then every $f_\alpha$ is holomorphic; hence $F\in\A(U)$ and has the canonical halo evaluation.
\end{proposition}
\begin{proof}
Uniqueness of Hahn coefficients gives $\int_{\partial T}f_\alpha\dd z=0$ for every coefficient and triangle. On an ordinary disk, these identities make the integral from a fixed point path independent along polygonal paths. The resulting primitive is complex differentiable, with derivative $f_\alpha$ by continuity. Its derivative is holomorphic, equivalently the usual scalar Morera theorem applies. Thus each coefficient is holomorphic on $U$.
\end{proof}
Continuity here is coefficientwise ordinary continuity. It is not a criterion for all fine-continuous surcomplex functions.

\section{Hahn--Weierstrass preparation and division}\label{sec:preparation}
This section supplies the main mechanism behind the global zero and residue theorems. The crucial point is that the perturbation may involve an arbitrary well-ordered set of positive surreal exponents, not merely integral powers of one small parameter.

\subsection{An ordinary division operator}
Let $V\subseteq\C$ be a domain containing all the roots of a monic polynomial $P_0(z)\in\C[z]$ of degree $N\geq1$. There are linear maps
\[
 T:\hol(V)\to\C[z]_{<N},\qquad
 Q_0:\hol(V)\to\hol(V)
\]
satisfying
\begin{equation}\label{eq:ordinarydivision}
 f=T(f)+P_0Q_0(f).
\end{equation}
To construct $T(f)$, prescribe at each root $a$ of $P_0$ the derivatives of $f$ through order one less than its multiplicity. Hermite interpolation gives the unique polynomial of degree less than $N$ with those data. Then $(f-T(f))/P_0$ has removable singularities at the roots and defines $Q_0(f)$. The decomposition is unique: a degree-less-than-$N$ polynomial divisible by $P_0$ is zero. Both operators extend coefficientwise to Hahn sections without increasing their supports.

\subsection{Preparation by well-founded recursion}
\begin{theorem}[Hahn--Weierstrass preparation]\label{thm:preparation}
Let
\[
 H(z)=P_0(z)+E(z),\qquad E\in\A_+(V).
\]
There exist unique factors
\begin{align*}
 P(z)&=P_0(z)+p(z), &\deg p&<N, &p&\in\infinit[z],\\
 W&=1+U, && &U&\in\A_+(V),
\end{align*}
such that $H=PW$. Here $p\in\infinit[z]$ means that every coefficient of $p$ is infinitesimal.

More precisely, if $S\subseteq\No_{>0}$ is the support of $E$, then the supports of $p$ and $U$ lie in $\langle S\rangle\setminus\{0\}$. The factor $W$ is a unit in $\A(V)$ and $\ev W$ is nowhere zero on $\halo V$.
\end{theorem}
\begin{proof}
Write the desired identity as
\begin{equation}\label{eq:prepidentity}
 E=p+P_0U+pU.
\end{equation}
Let $M=\langle S\rangle$, a well-ordered set by Lemma~\ref{lem:neumann}, and let $e_\delta$ denote the coefficient of $E$ at $\delta$, zero when $\delta\notin S$. Recursively in $\delta\in M_{>0}$ define
\begin{align}
 r_\delta&=e_\delta-
       \sum_{\substack{\alpha,\beta\in M_{>0}\\\alpha+\beta=\delta}}
              p_\alpha u_\beta,\label{eq:preprecursion1}\\
 p_\delta&=T(r_\delta),\qquad
 u_\delta=Q_0(r_\delta).\label{eq:preprecursion2}
\end{align}
The sum is finite. Its indices satisfy $\alpha,\beta<\delta$, so all its entries have already been constructed. Thus ordinary transfinite recursion on the well-ordered set $M_{>0}$ defines all coefficients, including at limit stages without taking analytic limits.

Each $p_\delta$ is a polynomial of degree less than $N$, and each $u_\delta$ is holomorphic on the same $V$. The series
$p=\sum_{\delta>0}t^\delta p_\delta$ and
$U=\sum_{\delta>0}t^\delta u_\delta$ are therefore Hahn sections with positive support. Since the polynomial degrees are uniformly bounded, the coefficients of $p$ are actual elements of $\infinit$, giving a single polynomial over $\K$. Equation~\eqref{eq:ordinarydivision} makes \eqref{eq:prepidentity} hold at every exponent.

For uniqueness, suppose $(P_0+p)(1+U)=(P_0+\widetilde p)(1+\widetilde U)$ with the stated positive-support conditions. If either difference is nonzero, let $\delta$ be the least exponent occurring in their combined support. Terms involving a difference multiplied by a positive-support correction cannot contribute at $\delta$. The coefficient of the difference of the products is therefore
\[
 (p-\widetilde p)_\delta+P_0(U-\widetilde U)_\delta=0.
\]
Uniqueness of \eqref{eq:ordinarydivision} forces both coefficients to vanish, a contradiction. Finally, $1+U$ has coherent inverse $\sum_{n\geq0}(-U)^n$ by the support lemma, and its evaluated value is one plus an infinitesimal at every halo point.
\end{proof}

\begin{remark}[Why ordinary completeness is not being used]
In a high-rank Hahn field, a sequence of corrections of valuations $n\delta$ need not become smaller than every surreal tolerance. A metric contraction proof using those corrections would therefore require justification not available from ordinary completeness. Equations~\eqref{eq:preprecursion1}--\eqref{eq:preprecursion2} avoid that issue: coefficients are assigned on a well-ordered support, with finitely many earlier contributions at each exponent. The construction is a Hahn sum, not a fine limit of successive approximations.
\end{remark}

\begin{theorem}[Hahn--Weierstrass division]\label{thm:division}
Let $P=P_0+p$ be as in Theorem~\ref{thm:preparation}. Every $A\in\A(V)$ has a unique decomposition
\begin{equation}\label{eq:Hahndivision}
 A=PQ+R,
 \qquad Q\in\A(V),\quad R\in\K[z],\quad\deg R<N.
\end{equation}
\end{theorem}
\begin{proof}
Extend $T,Q_0$ coefficientwise and put $\mathcal L(Q)=Q_0(pQ)$. Solving for the quotient requires
\[
 Q=Q_0(A-pQ),\qquad
 Q=\sum_{n\geq0}(-\mathcal L)^nQ_0(A).
\]
This operator geometric series is well defined: each application of $\mathcal L$ adds an exponent from the positive support of $p$, while $Q_0$ changes only coefficient functions. If $S_A$ is the support of $A$ and $S_p$ the positive support of $p$, the full family is supported in $S_A+\langle S_p\rangle$ with finite contributions at each exponent. Thus $Q\in\A(V)$ and $(I+\mathcal L)Q=Q_0(A)$.

Set $R=T(A-pQ)$. Applying \eqref{eq:ordinarydivision} to $A-pQ$ gives $A-pQ=P_0Q+R$, proving \eqref{eq:Hahndivision}. The remainder has uniformly bounded degree, so it is a polynomial over $\K$.

For uniqueness take the least Hahn exponent at which two proposed quotients or remainders differ. Since $p$ has positive support, the leading equation is $P_0q+r=0$, with $q\in\hol(V)$ and $\deg r<N$. Ordinary uniqueness forces both to be zero, a contradiction.
\end{proof}

\subsection{Preparation on a bounded complex domain}
\begin{corollary}[Finite-domain preparation]\label{cor:domainprep}
Let $D$ be a bounded Jordan domain with piecewise $C^1$ boundary, $\overline D\subseteq U$, and let
\[
 F=t^\beta(f_0+E)\in\A(U),\quad E\in\A_+(U),\quad
 f_0\not\equiv0,
\]
with $f_0$ having no zero on $\partial D$. Let $P_0$ be the monic complex polynomial consisting of all zeros of $f_0$ in $D$, with multiplicities, and let $N=\deg P_0$. There is a neighborhood $V$ of $\overline D$ on which
\begin{equation}\label{eq:domainprep}
 F=t^\beta u_0PW,
 \qquad u_0=f_0/P_0\in\hol(V)^\times,
\end{equation}
where $P=P_0+p$ is monic of degree $N$ with infinitesimal coefficient correction and $W\in1+\A_+(V)$ is a unit. When $N=0$, take $P=1$ and $W=1+E/f_0$.
\end{corollary}
\begin{proof}
There are finitely many zeros in $\overline D$ because the zeros of the nonzero ordinary holomorphic function $f_0$ are isolated and $\overline D$ is compact. Choose a neighborhood $V\subseteq U$ of $\overline D$ with no further zeros of $f_0$. Division by $P_0$ gives a nowhere-zero holomorphic $u_0$ on $V$. Apply Theorem~\ref{thm:preparation} to
$t^{-\beta}F/u_0=P_0+E/u_0$. The case $N=0$ is immediate.
\end{proof}
Compactness has been used only on an ordinary complex compact set, not on a surreal ball.

\section{Exact zero lifting and Rouch\'e stability}\label{sec:zeros}

\begin{theorem}[Zero clusters with exact multiplicity]\label{thm:zerolift}
In the setting of Corollary~\ref{cor:domainprep}, $\ev F$ has exactly $N$ zeros in $\halo D$, counted with local analytic multiplicity. If $a\in D$ is a zero of $f_0$ of multiplicity $m$, then precisely $m$ of these zeros have standard part $a$. There are no zeros whose standard parts are not zeros of $f_0$.
\end{theorem}
\begin{proof}
All factors in \eqref{eq:domainprep} except $P$ evaluate to nonzero functions on $\halo V$. Hence the zeros and their local multiplicities are exactly those of $P$.

The polynomial $P$ splits over $\K$, say $P(z)=\prod_{j=1}^N(z-\lambda_j)$, listing roots with multiplicity. Its coefficients are finite, because they are complex coefficients plus infinitesimals. No root can be infinite. Indeed, if $\eta=\vH(\lambda)<0$, the monic term $\lambda^N$ has valuation $N\eta$, while each lower term has valuation at least $k\eta>N\eta$, and cannot cancel it.

Every $\lambda_j$ therefore has a standard part. Applying the standard-part ring homomorphism to the factorization gives
\[
 P_0(z)=\prod_{j=1}^N(z-\st\lambda_j)
 \quad\text{in }\C[z].
\]
This proves the exact distribution of multiplicities among the ordinary roots of $P_0$. All of those roots lie in $D$. Away from zeros of $f_0$, equation~\eqref{eq:leadingcoherent} already excludes a zero. Multiplication by the analytic units in \eqref{eq:domainprep} does not change local multiplicities.
\end{proof}

\begin{corollary}[Simple-root lifting]
If $f_0(a)=0$ and $f_0'(a)\ne0$, then $\ev F$ has a unique zero in $a+\infinit$, and it is simple. Its displacement can be determined recursively as a Hahn series from the coefficient functions of $F$.
\end{corollary}
For example, when $F=f_0+t^\rho g+\text{terms of larger Hahn exponent}$ with $\rho>0$, the lifted simple root satisfies
\[
 \lambda=a-t^\rho\frac{g(a)}{f_0'(a)}+o(t^\rho).
\]
To see the leading correction directly, Taylor-expand at $a$. The linear term $f_0'(a)(\lambda-a)$ must cancel the first nonzero perturbation; the quadratic and higher terms have larger valuation. If $g(a)=0$, this displayed first correction is zero and the actual displacement is smaller.

\begin{corollary}[Infinitesimal Rouch\'e theorem]\label{cor:infrouche}
Suppose $F=t^\beta(f_0+E)$ satisfies the hypotheses of Theorem~\ref{thm:zerolift}, and $G\in\A(U)$ has $\nu(G)>\beta$. Then $\ev F$ and $\ev{F+G}$ have the same number of zeros in $\halo D$, counted with multiplicity. They have the same multiplicity total in every individual halo of an ordinary zero of $f_0$.
\end{corollary}
\begin{proof}
The two sections have the same leading ordinary coefficient $f_0$. Apply Theorem~\ref{thm:zerolift} to each.
\end{proof}

\begin{theorem}[Rouch\'e theorem with ordinary leading comparison]\label{thm:Rouche}
Let $F,G\in\A(U)$ and suppose, for some $\beta\in\No$,
\[
 t^{-\beta}F=f+E_F,\qquad
 t^{-\beta}G=g+E_G,\qquad E_F,E_G\in\A_+(U),
\]
where $f,g\in\hol(U)$ and
$|g(z)|<|f(z)|$ for $z\in\partial D$. Then $\ev F$ and $\ev{F+G}$ have the same number of zeros in $\halo D$, counted with multiplicity.
\end{theorem}
\begin{proof}
The boundary inequality makes both $f$ and $f+g$ nonzero on $\partial D$. Ordinary Rouch\'e gives the same total multiplicity for their zeros in $D$. Theorem~\ref{thm:zerolift} transfers those respective totals exactly to $F$ and $F+G$.
\end{proof}
The Hahn perturbations in these statements need not be uniformly bounded by one ordinary analytic majorant. Their common positive well-ordered support is what makes the preparation and zero count possible.

\subsection{Two examples of zero splitting}
\begin{example}[Scales beyond all integral powers]\label{ex:twoscale}
Let $t=\omega^{-1}$ and consider
\[
 F(z)=z^2-t-t^\omega z.
\]
Its leading ordinary coefficient is $z^2$. The two roots are
\[
 \lambda_\pm=\frac{t^\omega}{2}
 \pm t^{1/2}\left(1+\frac{t^{2\omega-1}}4\right)^{1/2},
\]
with the square root expanded at one. Both have standard part zero. The correction $t^\omega$ is smaller than every $t^n$ with standard integer $n$, but is retained exactly by the theory. The root count is two regardless of that scale separation.
\end{example}

\begin{example}[Two finite roots and one infinite root]\label{ex:cubic}
Consider
\[
 G(z)=z^2-t(1+z^3),\qquad t=\omega^{-1}.
\]
The ordinary leading coefficient is $z^2$. On any fixed sufficiently small standard disk about zero, Theorem~\ref{thm:zerolift} gives two roots, not three. With $q=t^{1/2}$ their first terms are
\begin{align*}
 \lambda_+&=q+\tfrac12q^4+\tfrac58q^7+q^{10}
                   +\tfrac{231}{128}q^{13}+\cdots,\\
 \lambda_-&=-q+\tfrac12q^4-\tfrac58q^7+q^{10}
                   -\tfrac{231}{128}q^{13}+\cdots.
\end{align*}
The third polynomial root is infinite. Vieta's formula gives
$\lambda_\infty=t^{-1}-(\lambda_++\lambda_-)$.

The preparation can be made explicit. Let $u=1-t^3u^{-2}$ be the unique formal series with constant term one. Then
\[
 u=1-t^3-2t^6-7t^9-30t^{12}-\cdots,
\]
and
\begin{equation}\label{eq:cubicfactor}
 G(z)=
 \left(z^2-\frac{t^2}{u^2}z-\frac tu\right)(u-tz).
\end{equation}
The first factor is the degree-two prepared polynomial; the second is a coherent unit on every bounded standard domain's halo. Its zero $u/t$ is infinite and is outside those halos. Identity~\eqref{eq:cubicfactor} follows by multiplication and $u=1-t^3/u^2$. The accompanying symbolic checks verify the displayed finite expansions, while the recursion and the preceding theorems justify the full Hahn identities.
\end{example}

\section{Moving poles, the residue theorem, and the argument principle}\label{sec:residuetheorem}

\subsection{Coherent meromorphic quotients}
A \emph{coherent meromorphic quotient} on $\halo U$ is represented by $H=F/G$, with $F,G\in\A(U)$ and $G\ne0$. Its evaluated values are $\ev F/\ev G$ wherever the denominator is nonzero, with cancellations interpreted as meromorphic germs. Equivalent fractions have the same values and germs because $\A(U)$ is a domain and evaluation is faithful.

Let $D$ be as in the preceding section, and assume the leading ordinary coefficient of $G$ has no zeros on $\partial D$. On a sufficiently narrow ordinary collar $C$ of that boundary, Proposition~\ref{prop:unit} makes $G^{-1}$ a coherent section. Hence $H$ has an unambiguous coherent expansion on $C$, and its coefficientwise contour integral is defined there.

\begin{theorem}[Holomorphic-plus-rational decomposition]\label{thm:meromdecomp}
Under these hypotheses, after choosing a neighborhood $V$ of $\overline D$ as in Corollary~\ref{cor:domainprep}, the quotient has a decomposition
\begin{equation}\label{eq:meromdecomp}
 H(z)=Q(z)+\frac{R(z)}{P(z)},
 \qquad Q\in\A(V),\quad R,P\in\K[z],\quad\deg R<\deg P,
\end{equation}
where the roots of $P$ all lie in $\halo D$. When the prepared denominator has degree zero, the rational term is absent. There are only finitely many genuine poles in $\halo D$.
\end{theorem}
\begin{proof}
Prepare the denominator as $G=t^\beta u_0PW$. The section
$A=t^{-\beta}F/(u_0W)$ is coherent and holomorphic on $V$, by Proposition~\ref{prop:unit}. Apply Hahn--Weierstrass division to $A$, obtaining $A=PQ+R$. This gives \eqref{eq:meromdecomp}. All roots of $P$ are in $\halo D$ by Theorem~\ref{thm:zerolift}, and cancellation can only remove poles from this finite list.
\end{proof}
This decomposition is considerably stronger than a coefficientwise residue calculation at fixed complex poles. It separates the actual poles in $\K$, including infinitesimally split ones, from a coherent holomorphic remainder.

\begin{theorem}[Surcomplex residue theorem]\label{thm:residue}
Let $H=F/G$ satisfy the preceding hypotheses, and let $p_1,\dots,p_s$ be its distinct genuine poles in $\halo D$. Then
\begin{equation}\label{eq:residuetheorem}
 \boxed{\quad
 \hoint_{\partial D}H(z)\dd z
 =2\pi i\sum_{j=1}^s\Res_{p_j}(\ev H(z)\dd z).
 \quad}
\end{equation}
Each residue is defined in the local surcomplex coordinate $z-p_j$. The sum is finite, and its terms may be infinite surcomplex numbers.
\end{theorem}
\begin{proof}
The coherent holomorphic term $Q$ in \eqref{eq:meromdecomp} integrates to zero by \eqref{eq:Cauchytheorem}. Over the algebraically closed field $\K$, the proper rational term has a finite partial-fraction decomposition
\[
 \frac RP=\sum_{j=1}^s\sum_{k=1}^{m_j}
                  \frac{c_{jk}}{(z-p_j)^k}.
\]
Its local residue at $p_j$ is $c_{j1}$. Write $p_j=a_j+\epsilon_j$, with $a_j\in D$ and $\epsilon_j\in\infinit$. On the boundary collar,
\[
 (z-p_j)^{-k}=\sum_{n\geq0}
 \binom{n+k-1}{k-1}\epsilon_j^n(z-a_j)^{-n-k}.
\]
This expansion is Hahn-summable. Ordinary contour integration is zero for every displayed term except $k=1,n=0$, whose integral is $2\pi i$. Multiplication by $c_{jk}$ and finite addition now give exactly \eqref{eq:residuetheorem}.
\end{proof}
The proof avoids replacing distinct nearby surcomplex poles by their standard parts when computing local residues. Standard parts organize the clusters; the residues themselves remain surcomplex.

\begin{theorem}[Argument principle]\label{thm:argument}
Let $F,G\in\A(U)$ be nonzero, with both leading ordinary coefficients nonzero on $\partial D$, and let $H=F/G$. Then
\begin{equation}\label{eq:argument}
 \frac1{2\pi i}\hoint_{\partial D}\frac{H'(z)}{H(z)}\dd z
 =\sum_{p\in\halo D}\ord_p\ev H
 =N_F-N_G\in\Z,
\end{equation}
where $N_F,N_G$ count zeros with multiplicity in $\halo D$. Only finitely many terms on the middle side are nonzero, and cancellations are automatically counted with their signs.
\end{theorem}
\begin{proof}
The logarithmic derivative is a coherent meromorphic quotient on the boundary collar. Its poles in $\halo D$ occur only at zeros and poles of $H$. By Theorem~\ref{thm:formalprimitive}, its residue there is $\ord_pH$. Apply Theorem~\ref{thm:residue}. At every point, $\ord_p(F/G)=\ord_pF-\ord_pG$, so the sum equals $N_F-N_G$.
\end{proof}
The normalized integral is an exact ordinary integer, even though the individual zeros, derivatives, and residues used in its derivation live in $\K$. It can be used as an analytic index. We do not identify it with the winding number of a continuous loop in the fine topology, where such an ordinary nonconstant loop is unavailable.

\subsection{Primitives and removable meromorphic singularities}
\begin{corollary}[Residues are the finite-domain primitive obstruction]\label{cor:resprimitive}
On the halo of the Jordan domain $D$, a coherent meromorphic quotient as above has a coherent meromorphic primitive if and only if all its local residues vanish.
\end{corollary}
\begin{proof}
Necessity is local. For sufficiency, use \eqref{eq:meromdecomp}. The coherent holomorphic term has a primitive on the simply connected domain $D$ by Theorem~\ref{thm:coherentprimitive}. In the finite partial-fraction part, absence of residues removes every $k=1$ term. Integrating the terms with $k\geq2$ gives rational functions with surcomplex coefficients. Their sum and the coherent primitive form a coherent meromorphic quotient.
\end{proof}
Equivalently, for any fixed finite permitted set of poles, meromorphic differentials modulo derivatives are classified by their residue vectors in a finite-dimensional $\K$-vector space. Surjectivity is witnessed by the differentials $\dd z/(z-p)$.

\begin{corollary}[Removability within the meromorphic class]\label{cor:removable}
A bounded meromorphic germ on a punctured fine neighborhood has no pole and hence extends analytically. If a coherent quotient in Theorem~\ref{thm:meromdecomp} has no genuine poles in $\halo D$, it extends as a coherent holomorphic section on a neighborhood of $\overline D$.
\end{corollary}
\begin{proof}
A pole of order $m>0$ has the form $c h^{-m}(1+\theta(h))$, with $c\ne0$ and $\theta(h)\to0$ finely. Given any fixed bound, choose a sufficiently small nonzero $h$ so that this expression exceeds the bound in modulus. This contradicts boundedness. For the global statement, a proper rational function $R/P$ with no poles over the algebraically closed field $\K$ is zero: after cancellation its denominator is constant, but a proper rational function cannot be a nonzero polynomial. Thus only the coherent holomorphic term $Q$ remains.
\end{proof}
The first assertion assumes a meromorphic germ, not an arbitrary analytic function on a punctured neighborhood. That distinction is essential; a counterexample appears in Section~\ref{sec:limits}.

\subsection{A residue calculation with infinite summands}
\begin{example}[Two poles coalescing at an ordinary point]\label{ex:residuesplit}
Let $t=\omega^{-1}$ and let the standard unit circle be positively oriented. Consider
\[
 H(z)=\frac{\ev{\exp}(z)}{z^2-t}.
\]
The poles in the finite halo are $p_\pm=\pm\sqrt t$, with residues
\[
 r_+=\frac{\ev{\exp}(\sqrt t)}{2\sqrt t},\qquad
 r_-=-\frac{\ev{\exp}(-\sqrt t)}{2\sqrt t}.
\]
Each residue is infinite in modulus, but their sum is finite:
\[
 r_++r_-=
 \frac{\ev{\sinh}(\sqrt t)}{\sqrt t}
 =\sum_{n\geq0}\frac{t^n}{(2n+1)!}.
\]
Therefore
\begin{equation}\label{eq:residueexample}
 \hoint_{|z|=1}\frac{e^z}{z^2-t}\dd z
 =2\pi i\sum_{n\geq0}\frac{t^n}{(2n+1)!}.
\end{equation}
In the contour integrand $e^z$ denotes its ordinary coefficient function; its halo evaluation is $\ev{\exp}$. Directly expanding
$(z^2-t)^{-1}=\sum_{n\geq0}t^nz^{-2n-2}$ on the contour gives the same answer by ordinary coefficientwise residues. Thus the calculation compares two genuinely different organizations of the same exact result: split poles with half-integral scales, and an integral-power Hahn expansion at the ordinary center.
\end{example}

\section{Local differential equations}\label{sec:ODE}
The universal realization theorem has a useful consequence beyond the usual one-variable complex-analytic list.

\begin{theorem}[Analytic ODEs have unique analytic solution germs]\label{thm:ODE}
Let $A(X,Y)$ be a vector of Hahn-analytic germs near $(0,b)\in\K\times\K^q$. The initial value problem
\[
 Y'(X)=A(X,Y(X)),\qquad Y(0)=b
\]
has a unique Hahn-analytic solution germ.
\end{theorem}
\begin{proof}
Replace $Y$ by $b+Z$ so that the unknown $Z$ has zero constant term. Write $Z(X)=\sum_{n\geq1}z_nX^n$. At degree $n$, the equation specifies
\[
 (n+1)z_{n+1}=[X^n]A\left(X,b+\sum_{j=1}^{n}z_jX^j\right).
\]
Only finitely many formal coefficient expressions contribute at a fixed degree because $Z(0)=0$. Thus the recursion determines every $z_n$ uniquely. The resulting coefficient sequence is a set, so the formal solution is locally realized by Theorem~\ref{thm:realization}. The several-variable composition and derivative theorems realize its equation. Any analytic solution germ has the same formal coefficients by the recursion.
\end{proof}
No numerical majorant estimates are needed: a formal solution with very rapidly growing coefficients still has a sufficiently small Hahn-analytic domain. This is stronger than merely transferring the classical proof, but weaker than a global existence statement. The allowed neighborhood may be much smaller than any prescribed standard or previously chosen infinitesimal scale. Uniqueness is uniqueness of a germ; disconnected fine domains still permit unrelated continuation elsewhere.

For linear systems $Y'=A(X)Y+B(X)$, the same recursion constructs a fundamental matrix with any prescribed initial matrix. An initially invertible fundamental matrix stays invertible on a sufficiently small neighborhood by the local inverse-of-a-unit argument. Formal identities such as variation of constants are therefore valid locally after realization.

\section{Liouville theorems, growth, and precise limitations}\label{sec:limits}

\subsection{Why boundedness alone does not suffice}
Example~\ref{ex:badliouville} already gives a bounded nonconstant Hahn-analytic function on all of $\K$. Restricting to coherent sections on a finite halo does not restore a naive Liouville theorem either.

\begin{example}[A real bound can hide an unbounded coefficient]\label{ex:tz}
The coherent entire section $F(z)=tz$, $t=\omega^{-1}$, satisfies
\[
 |\ev F(z)|<1\qquad\text{for every }z\in\halo\C=\fin,
\]
yet is nonconstant. Each individual $z$ is finite, so multiplying it by $t$ is infinitesimal. The coefficient function $z$ is unbounded on $\C$, but its size is hidden behind the infinitesimal scale $t$.
\end{example}
Even the section $F(z)=z$ is bounded on $\fin$ by the infinite constant $\omega$. Thus there are two distinct pitfalls: a finite-halo domain omits infinite arguments, and ordered-field inequalities do not resolve every Hahn coefficient.

\begin{theorem}[Coefficientwise polynomial-growth theorem]\label{thm:Liouville}
Let $F=\sum_\gamma t^\gamma f_\gamma\in\A(\C)$, and let $d\geq0$ be an integer. Suppose for each $\gamma$ there is a real constant $C_\gamma\geq0$ such that
\[
 |f_\gamma(z)|\leq C_\gamma(1+|z|)^d
 \qquad(z\in\C).
\]
Then $F$ is a polynomial of degree at most $d$ with coefficients in $\K$. In particular, if every coefficient function is bounded on $\C$, then $F$ is a constant surcomplex section. Conversely every polynomial of degree at most $d$ admits such coefficientwise bounds.
\end{theorem}
\begin{proof}
For each fixed $\gamma$, the ordinary Cauchy estimate at zero gives
\[
 \frac{|f_\gamma^{(n)}(0)|}{n!}
 \leq \frac{C_\gamma(1+R)^d}{R^n}
 \qquad(R>0\text{ real}).
\]
For $n>d$, let the ordinary real radius $R$ tend to infinity. This gives $f_\gamma^{(n)}(0)=0$, so the entire function $f_\gamma$ is a polynomial of degree at most $d$. Collecting its finitely many powers gives
$F(z)=\sum_{j=0}^d a_jz^j$ with
$a_j=\sum_\gamma t^\gamma f_\gamma^{(j)}(0)/j!\in\K$.
For the converse expand each of the finitely many $a_j$ in normal form, take the finite union of their supports, and bound each coefficient polynomial by the sum of the moduli of its coefficients times $(1+|z|)^d$.
\end{proof}
The hypothesis controls each scale. It is not equivalent to the existence of a single surreal number bounding $\ev F$.

\begin{proposition}[Algebraic Liouville theorem on the full field]\label{prop:rationalLiouville}
If a rational function $R\in\K(z)$ has no poles anywhere in $\K$ and is bounded on all of $\K$ by a fixed surreal constant, then $R$ is constant.
\end{proposition}
\begin{proof}
Write $R=P/Q$ in reduced form. If $Q$ were nonconstant, algebraic closedness would provide a root at which $R$ has a pole. Thus $R$ is a polynomial. For a nonconstant polynomial $a_nz^n+\cdots+a_0$, take a positive surreal $z$ sufficiently large that the leading term exceeds twice the sum of the moduli of all other terms and also exceeds twice the asserted bound. The reverse triangle inequality then contradicts boundedness. Only finitely many inequalities are being imposed.
\end{proof}
This proposition uses all surcomplex arguments, unlike Theorem~\ref{thm:Liouville}, whose global domain is a finite halo. It supplies a genuine full-field statement without suggesting that the analytic proof of classical Liouville works unmodified.

\subsection{A bounded punctured analytic function need not be removable}
\begin{example}[Failure of unrestricted Riemann removability]\label{ex:nonremovable}
For $z\ne0$, let $\lc(z)\in\C^\times$ be its leading Hahn coefficient and define
\[
 B(z)=\begin{cases}
 1,&\operatorname{Re}\lc(z)>0,\\
 0,&\operatorname{Re}\lc(z)\leq0.
 \end{cases}
\]
If $h\prec z$, then $z+h$ has the same leading exponent and coefficient as $z$. Therefore $B$ is locally constant at every nonzero point: the ball with radius $|z|/\omega$ suffices. It is Hahn analytic on every punctured fine neighborhood of zero and bounded by one.

Nevertheless $B$ has no limit at zero. Every neighborhood contains a positive monomial $r$ and its negative, on which the values are respectively one and zero. An assumed limit would require both values to lie within, for example, $1/3$ of that limit, contradicting the triangle inequality. Thus no continuous, let alone analytic, extension at zero exists.
\end{example}
This example shows why Corollary~\ref{cor:removable} was explicitly restricted to meromorphic germs and coherent meromorphic quotients.

\subsection{Other assertions deliberately not made}
There is no general radius-of-convergence formula based on ordinary complex coefficient magnitudes for arbitrary surcomplex coefficients. A Hahn sum is not a fine-topological limit of its ordinary partial sums. Agreement on ordinary points does not determine an arbitrary fine-analytic function. A real-parameter contour is not an ordinary continuous path in the fine topology. A boundedness hypothesis in the ordered field does not automatically bound every coefficient.

Nor have we constructed a canonical single-valued exponential on all of $\K$, an unrestricted global contour integral for arbitrary class functions, a universal analytic continuation across all scales, or a general theory of essential singularities with a Picard theorem. These are separate structural questions. The explicit theorems above remain valid without resolving them or silently assuming them.

\section{Affine charts at infinitesimal and infinite scales}\label{sec:charts}
The finite-halo construction is not confined to neighborhoods of finite points. It can be transported to any center and any nonzero scale.

\begin{definition}[Scaled coherent chart]
For $b\in\K$, $r\in\K^\times$, and an ordinary complex domain $U$, put
\[
 \mathcal U(b,r;U)=\{b+r\xi:\xi\in\halo U\}.
\]
A function on this chart is \emph{coherent in the coordinate $(z-b)/r$} if it has the form
\[
 f(z)=\ev F\bigl((z-b)/r\bigr),\qquad F\in\A(U).
\]
The affine coordinate is part of the specified structure.
\end{definition}
Every such chart is fine open. The center $b$ may have infinite real or imaginary part; the modulus of $r$ may be infinitesimal, finite, or infinite.

\begin{theorem}[Affine transport]\label{thm:affine}
In a scaled coherent chart, the identity, open-mapping, local maximum-modulus, preparation, and zero-counting theorems transport from the corresponding theorems on $\halo U$. Derivatives satisfy
\[
 f^{(n)}(b+r\xi)=r^{-n}\ev{D^nF}(\xi).
\]
For a standard piecewise $C^1$ path $\gamma$ in $U$, define the scaled contour functional by
\begin{equation}\label{eq:scaledcontour}
 \hint_{b+r\gamma} f(z)\dd z
 :=r\hint_\gamma F(\xi)\dd\xi.
\end{equation}
Cauchy's formula and the residue theorem hold with this functional. If $f$ is meromorphic and $p=b+r\lambda$, then
\begin{equation}\label{eq:scaledresidue}
 \Res_{p}\bigl(f(z)\dd z\bigr)
 =r\Res_{\lambda}\bigl(\ev F(\xi)\dd\xi\bigr).
\end{equation}
In particular, the zero and pole multiplicities, and the normalized logarithmic-derivative index, are unchanged by the affine coordinate.
\end{theorem}
\begin{proof}
The affine map $\xi\mapsto b+r\xi$ is a fine homeomorphism, with analytic inverse. The derivative formula follows by the chain rule; local multiplicities are preserved by its nonzero linear term. Pulling a preparation factorization back gives the factorization in the chart coordinate, with the same finite number of roots. Equation~\eqref{eq:scaledcontour} is the definition of the transported functional, so its Cauchy and residue formulas follow by substitution.

For the explicit residue calculation, write $F(\lambda+X)=\sum_n a_nX^n$. Then
$f(p+Y)=\sum_n a_nr^{-n}Y^n$, and its $Y^{-1}$ coefficient is $ra_{-1}$, proving \eqref{eq:scaledresidue}. In the logarithmic derivative, the factor $r^{-1}$ from differentiation cancels the factor $r$ from the contour differential. The resulting index is therefore exactly the one in Theorem~\ref{thm:argument}.
\end{proof}
For example, one may take $b=\omega+i\omega^2$ and $r=\omega^{-\omega}$. The resulting domain is an extremely small neighborhood of an infinite point, and its zero clusters and residues are just as accessible as those in a standard unit disk. Taking $r=\omega^\omega$ gives a chart at an infinitely large scale instead.

The notation $b+r\gamma$ in \eqref{eq:scaledcontour} records an affine contour datum, not a nonconstant ordinary continuous path into the fine topology. This caveat does not disappear when the center or scale changes.

Ordinary complex affine coordinate changes preserve the coherent class by coefficientwise composition. Arbitrary changes involving surreal centers and scales may relate very different ordinary base domains; no unrestricted gluing or chart-independence theorem is asserted. An intrinsic global theory would need to specify admissible transitions and prove their compatibility with integration.

\section{What the two theories include, and what connects them}\label{sec:comparison}

\subsection{The coherent class is genuinely smaller than the local class}
The additional global structure is not merely terminology.

\begin{proposition}[A local germ with no coherent extension]\label{prop:strictness}
The germ at zero of
\[
 E(X)=\sum_{n\geq0}n!X^n
\]
is Hahn analytic, but there is no connected ordinary neighborhood $U$ of zero and no $F\in\A(U)$ whose evaluated germ at zero is this germ.
\end{proposition}
\begin{proof}
Local analyticity was established in Example~\ref{ex:euler}. If such $F=\sum_\gamma t^\gamma f_\gamma$ existed, equality of analytic germs would give
\[
 \sum_\gamma t^\gamma\frac{f_\gamma^{(n)}(0)}{n!}=n!
 \qquad(n\geq0).
\]
Uniqueness of Hahn coefficients makes every derivative of $f_\gamma$ at zero vanish for $\gamma\ne0$. Ordinary analyticity and connectedness make each such coefficient function identically zero. The remaining coefficient would have ordinary Taylor series $\sum n!X^n$, whose radius is zero. This is impossible for a holomorphic function on a neighborhood of zero.
\end{proof}
Thus all formal germs can be realized locally, but they cannot all be placed in one fixed standard-coordinate coherent class. Affine rescaling can change this comparison: for example $E(tz)=\sum n!t^nz^n$ is a coherent section in the coordinate $z$. Scale is part of the analytic information, not a harmless normalization that can always be suppressed.

\subsection{Set-sized fields versus the full surreal field}
A calculation with specified coefficients and supports takes place inside some set-sized Hahn field, possibly after enlarging its exponent group. This makes coefficient recursion and exact symbolic work ordinary set-sized mathematics. The full proper-class field supplies additional scales smaller than every scale in that chosen field.

One must not transfer Proposition~\ref{prop:setdiscrete} to a fixed set-sized Hahn field with only its internal radii. The proof uses a positive surreal lower bound that need not belong to that field. Similarly, an arbitrary formal series over a fixed Hahn field may require adjoining smaller scales before the universal realization argument applies. The local realization theorem is a theorem about $\No[i]$, not an assertion that every fixed valued subfield is analytically complete in the same sense.

\subsection{A theorem map}
The principal conclusions can be organized by the extra structure they require. The entries below are statements proved in this article, with the classical and Hahn-field foundations specified in Section~1.

\begin{longtable}{@{}>{\raggedright\arraybackslash}p{0.28\textwidth}>{\raggedright\arraybackslash}p{0.40\textwidth}>{\raggedright\arraybackslash}p{0.26\textwidth}@{}}
\toprule
\textbf{Result} & \textbf{Required analytic structure} & \textbf{Reference}\\
\midrule
\endfirsthead
\toprule
\textbf{Result} & \textbf{Required analytic structure} & \textbf{Reference}\\
\midrule
\endhead
\bottomrule
\endfoot
Power-series realization & Arbitrary set of surcomplex coefficients; sufficiently small variable & Theorem~\ref{thm:realization}\\[5pt]
Cauchy--Riemann criterion & Total differentiability; real-analytic hypotheses for the converse to analyticity & Proposition~\ref{prop:CR}; Theorem~\ref{thm:analyticCR}\\[5pt]
Inverse, implicit, and branching theorems & Hahn-analytic germs & Theorems~\ref{thm:inverse}--\ref{thm:normalform}\\[5pt]
Formal residues and primitives & Meromorphic germs with finite principal parts & Theorem~\ref{thm:formalprimitive}\\[5pt]
Schwarz and Schwarz--Pick & Canonical lifts of ordinary disk maps & Theorems~\ref{thm:Schwarz}, \ref{thm:SchwarzPick}\\[5pt]
Identity and maximum modulus & Common-support coherent sections on a connected base & Theorem~\ref{thm:identity}; Corollary~\ref{cor:globalmapping}\\[5pt]
Cauchy and Morera & Coefficientwise contours and common support & Theorem~\ref{thm:Cauchy}; Proposition~\ref{prop:Morera}\\[5pt]
Preparation and division & Positive Hahn perturbation of a monic complex polynomial & Theorems~\ref{thm:preparation}, \ref{thm:division}\\[5pt]
Exact zero lifting and Rouch\'e & Leading ordinary coefficient nonzero on the boundary & Theorem~\ref{thm:zerolift}; Corollary~\ref{cor:infrouche}\\[5pt]
Residues and argument principle & Coherent meromorphic quotients with admissible boundary & Theorems~\ref{thm:residue}, \ref{thm:argument}\\[5pt]
Liouville-type conclusion & Coefficientwise growth, or a full-field rational function & Theorem~\ref{thm:Liouville}; Proposition~\ref{prop:rationalLiouville}\\[5pt]
Arbitrary centers and scales & A specified affine coherent chart & Theorem~\ref{thm:affine}\\
\end{longtable}

\section{Further development and concluding perspective}\label{sec:future}
The results suggest several concrete directions. These are proposed extensions of the present framework, not claims that the corresponding broad subjects have no existing literature.

\paragraph{Compatibility across scales.}
A useful next step is to characterize affine or nonlinear transition maps for which two coherent presentations induce the same analytic functions and contour functionals on overlaps. The main issue is not local differentiation, which already works, but preservation of a common well-ordered support over an ordinary base. A satisfactory atlas would have to retain sufficient information to avoid the locally constant extension ambiguity of Section~\ref{sec:topology}.

\paragraph{Several-variable coherent preparation.}
The local inverse and implicit theorems already work in finite dimension. A coherent several-variable preparation theorem would replace Hermite division by ordinary Weierstrass division with holomorphic parameter dependence. One would need a common parameter domain for all coefficient functions and support-preserving division operators. The positive-support recursion in Theorem~\ref{thm:preparation} identifies exactly where those inputs would enter.

\paragraph{Essential singularities and resummation.}
The local meromorphic theory uses finite principal parts. Allowing two-sided formal expansions requires new summability hypotheses: a support that is harmless in one variable may cease to be well ordered after evaluation. The factorial example also shows that local Hahn analyticity does not select a classical sectorial sum. Relating coherent charts to resurgent extensions such as those studied by Costin and Ehrlich \cite{CostinEhrlich} requires additional asymptotic data, not just equality of formal Taylor coefficients.

\paragraph{Larger contour structures.}
The present contour functional retains ordinary geometry and integrates coefficientwise. A broader theory might use explicitly controlled families of charts or cycles. It would have to state a summability rule for the entire family of integrals and prove invariance under admissible refinements. Neither ordinary real-parameter continuity nor an unnamed transfinite limit can substitute for those requirements.

\paragraph{Effective calculation and formal verification.}
For a finitely generated positive support, the preparation recursion gives an implementable exact algorithm to any specified finite support cutoff for which only finitely many terms occur. In higher rank, a bound on valuation need not leave only finitely many terms, so an algorithm must specify a finite support region or a finite truncation scheme rather than merely saying ``compute below precision $t^\delta$.'' The accompanying script checks several finite algebraic consequences. A machine-checked development would additionally formalize Hahn summability, normal forms, and the recursion theorem; the script is not a replacement for that task.

The overall conclusion is that surcomplex analysis is not simply complex analysis with larger constants. Its local side is in one respect more permissive: every formal power series is analytic on some sufficiently small scale. Its global side is more demanding: topology alone provides far too little analytic coherence. The common-support construction supplies one explicit bridge. With that bridge, local branching, Cauchy formulas, finite-domain preparation, moving poles, zero multiplicities, and residue indices fit into a single exact theory.

\appendix
\section{Notation and proof dependencies}
The symbol $t$ is the infinitesimal $\omega^{-1}$, and $t^\gamma$ always means the Conway monomial $\omega^{-\gamma}$. It is not being defined through a newly constructed surcomplex exponential. The valuation $v_H$ concerns Hahn support; the derivative $D$ concerns the analytic variable. The letters $X,Y$ are formal variables, whereas $z,w,\epsilon$ denote field elements when evaluated. A hat distinguishes an evaluated function from a formal coherent section whenever the distinction matters.

The notation $U^{\mathrm h}$ means $\st^{-1}(U)$ inside the finite elements of $\K$. It is not the whole fine ball obtained by replacing the complex coordinates of $U$ with arbitrary surreal coordinates. In particular, $\mathbb D^{\mathrm h}$ omits points infinitesimally close to the ordinary unit circle, even when their surcomplex modulus is strictly less than one.

The main proof chain is
\[
\begin{gathered}
 \text{normal forms and Hahn--Neumann support control}\\
 \Downarrow\\
 \text{local realization, differentiation, and analytic germs}\\
 \Downarrow\\
 \text{coherent evaluation and coefficientwise contour calculus}\\
 \Downarrow\\
 \text{well-founded preparation and division}\\
 \Downarrow\\
 \text{exact zero lifting and holomorphic-plus-rational decomposition}\\
 \Downarrow\\
 \text{surcomplex residues and the argument principle}.
\end{gathered}
\]
The algebraic closedness of $\K$ enters in extracting roots, local ramification, and splitting the prepared polynomial. Ordinary compactness enters only to obtain finitely many ordinary zeros in a bounded domain. Ordinary complex analysis is used only for coefficient functions or canonical lifts, never through an unstated transfer principle to all of $\K$.

\section{Reproducibility and status of the arguments}
The companion file \texttt{verify\_examples.py} uses exact rational symbolic arithmetic to check the displayed Catalan inverse, the factorial-series differential equation, the cubic's small-root expansions and preparation identity, and the cancellation formula for split-pole residues. Its output is supplied in \texttt{verification.txt}. These are consistency checks on finite truncations and algebraic identities. They do not check the class-theoretic foundations, replace the summability proofs, or constitute a proof-assistant certification of the article.

The literature discussion was checked against the cited primary sources and their available version metadata as of September 21, 2026. The article distinguishes established local surreal-analytic foundations from the explicit coherent framework developed here. No exhaustive novelty search or priority claim is intended.

\begin{thebibliography}{99}
\addcontentsline{toc}{section}{References}

\bibitem{Conway}
J.~H. Conway, \emph{On Numbers and Games}, Academic Press, London, 1976.

\bibitem{Gonshor}
H. Gonshor, \emph{An Introduction to the Theory of Surreal Numbers}, London Mathematical Society Lecture Note Series 110, Cambridge University Press, 1986.
\href{https://doi.org/10.1017/CBO9780511629143}{doi:10.1017/CBO9780511629143}.

\bibitem{Alling}
N.~L. Alling, \emph{Foundations of Analysis over Surreal Number Fields}, North-Holland Mathematics Studies 141, North-Holland, Amsterdam, 1987.

\bibitem{Neumann}
B.~H. Neumann, On ordered division rings, \emph{Transactions of the American Mathematical Society} \textbf{66} (1949).
\href{https://doi.org/10.1090/S0002-9947-1949-0032593-5}{doi:10.1090/S0002-9947-1949-0032593-5}.

\bibitem{BMgerms}
A. Berarducci and V. Mantova, Transseries as germs of surreal functions, \emph{Transactions of the American Mathematical Society} \textbf{371} (2019), 3549--3592.
\href{https://doi.org/10.1090/tran/7428}{doi:10.1090/tran/7428}.
\href{https://arxiv.org/abs/1703.01995v3}{arXiv:1703.01995v3}.

\bibitem{BMderiv}
A. Berarducci and V. Mantova, Surreal numbers, derivations and transseries, \emph{Journal of the European Mathematical Society} \textbf{20} (2018), 339--390.
\href{https://doi.org/10.4171/JEMS/769}{doi:10.4171/JEMS/769}.
\href{https://arxiv.org/abs/1503.00315}{arXiv:1503.00315}.

\bibitem{RSS}
S. Rubinstein-Salzedo and A. Swaminathan, Analysis on surreal numbers, \emph{Journal of Logic and Analysis} \textbf{6}, paper 5 (2014), 1--39.
\href{https://arxiv.org/abs/1307.7392v3}{arXiv:1307.7392v3}.

\bibitem{CostinEhrlich}
O. Costin and P. Ehrlich, Integration on the Surreals,
\href{https://arxiv.org/abs/2208.14331v5}{arXiv:2208.14331v5}, July 4, 2024.

\bibitem{MS}
J. M\"uller and A. Strohmaier, The theory of Hahn-meromorphic functions, a holomorphic Fredholm theorem, and its applications, \emph{Analysis \& PDE} \textbf{7} (2014), no.~3, 745--770.
\href{https://doi.org/10.2140/apde.2014.7.745}{doi:10.2140/apde.2014.7.745}.
\href{https://arxiv.org/abs/1205.0236v3}{arXiv:1205.0236v3}.

\bibitem{Mantova2026}
V. Mantova, Monotonicity and a Taylor approximation theorem for transseries,
\href{https://arxiv.org/abs/2601.07747v2}{arXiv:2601.07747v2}, May 9, 2026.

\end{thebibliography}
\end{document}
